\documentclass[11pt]{article}
\usepackage[margin=1.1in]{geometry}
\usepackage{amsmath,amssymb,amsthm}
\usepackage{graphicx}
\usepackage{booktabs}
\usepackage{hyperref}
\usepackage{microtype}
\usepackage{xcolor}
\usepackage{array}
\usepackage{enumitem}

\hypersetup{colorlinks=true, linkcolor=blue!60!black, urlcolor=blue!60!black}

\title{\textbf{VCML / VCSM: A Complete Theory of Viability-Gated Spatial Memory}\\[0.5em]
\large Theory Reference --- Papers 0--71}
\author{Gabriel Aubin-Moreau\\[0.3em]
\small VCSM Theoretical Series, March 2026}
\date{}

\begin{document}
\maketitle
\tableofcontents
\bigskip

% ─────────────────────────────────────────────────────────────────────────────
\begin{abstract}
We present the complete theory of the Viability-Constrained Coupled Map
Lattice (VCML) and its learning rule, the Viability-Gated Contrastive State
Memory (VCSM).
The central claim is that mandatory individual turnover is not an obstacle
to spatial memory but its primary mechanism.
Cells live, detect input, consolidate a contrastive memory trace, collapse
when unviable, and transmit their field state to successors via a copy-forward
birth step.
This local five-step rule produces: (I) spontaneous zone formation from
spatially patterned wave input; (II) analytically tractable parameter
boundaries for the adaptive regime; (III) a stochastic Allen-Cahn/Burgers
field equation in the spatial continuum limit; (IV) a mixed
pitchfork--spinodal commitment epoch at $t^* \approx 2000$ steps followed by
a metastability half-life of $\approx 800$ steps; (V) relay coupling between
regions that preserves zone structure across spatial gaps; (VI) continual
learning in which periodic task switching --- rather than eroding structure ---
amplifies it $9.7\times$ via copy-forward refreshing.
The system is formally non-gradient (curl $\neq 0$, machine-verified in Lean
4), lives in an 8D product manifold of viability and structure, and makes
quantitative predictions about hippocampal neurogenesis, adult memory
reconsolidation, and topographic map formation.
The theory's unique claim is about \emph{maintenance}, not formation:
zone formation is environmentally driven and robust to any consolidation rule;
zone maintenance under mandatory turnover requires viability-gated copy-forward.
Single-region theory constitutes a complete mechanistic account as of Paper~44:
Paper~42 closed all anomalies; Paper~43 demonstrated rule-family robustness,
confirming this distinction empirically (maintenance ratio $1.91\times$ with
viability gating vs.\ $1.11\times$ without; field-decay-only $= 0.64\times$).
Paper~44 closes the imprinting critique: under 2000 steps of sustained adversarial
input (wave-zone assignments permanently flipped), the adaptive mechanism maintains
positive fidelity with the original encoding ($+0.129$) while the passive mechanism
encodes the adversarial pattern ($-0.076$); $\Delta = 0.205$, $p < 0.001$.
Paper~45 validates the causal geometry conjecture externally:
ablation-influence partial correlations in GPT-2 flip sign at layer~4 from positive
(feature construction) to negative (causal integration, peak $r_p = -0.337$,
$p \ll 0.001$), with random-token control tracking ablation at every layer,
confirming that transformer geometry is organised by causal routing structure.
Paper~46 closes the temporal gate invariant and validates $\Omega = \mathrm{WR}\times\phi_w$
as the operative $C_2$ coordinate.
Paper~47 confirms $\delta = W_\mathrm{zone}/r_\mathrm{wave}$ as the third $C_2$ coordinate
(encoding quality peaks at $\delta=5$, degrades at $\delta \leq 2.5$) and reveals that
gate selectivity controls \emph{commitment rate} as well as adversarial filtering ---
the same mechanism at the same timescale.
Paper~48 establishes that VCML builds a \emph{population code}: single-cell zone decode
accuracy is only 4--5\% above chance across all gate conditions; zone information lives
in zone means, not individual cells.
The quiescence gate's operative function is noise reduction ($\sigma_w$ halves at high SS),
not amplitude enhancement; the true causal quality measure is
$\mathrm{SNR} = \mathrm{sg4}/\sigma_w$, which correctly ranks gate conditions.
The primitive underlying the entire architecture is A $\neq$ B (distinction detection);
viable/non-viable is one engineering instantiation.
The viability boundary is the maximum-information, maximum-distinction site in the
substrate; zone structure is the accumulated geometry of boundary crossings.
\end{abstract}

% ─────────────────────────────────────────────────────────────────────────────
\section{The System}

\subsection{The primitive: A $\neq$ B}

The ontological primitive of VCML is not viability.
It is \emph{distinction}: the detection of a difference between two states.

A cell that records the difference $h_i - \mathrm{baseline}_i$ is implementing A $\neq$ B.
It is asking: \emph{is the current state different from the expected state?}
The answer is a signed scalar that propagates through the three-timescale architecture
into the slow structural manifold $F_i$.
Everything else --- viability thresholds, collapse rates, zone geometry, the wave protocol ---
is engineering scaffolding that instantiates this primitive reliably.

Viable/non-viable is one instantiation of A $\neq$ B, chosen because it produces
a reliable, binary, threshold-crossing event (the collapse) that can serve as an
unambiguous write trigger.
It is not the primitive; it is the mechanism.
The primitive fires whenever the current hidden state differs from baseline.
The viability boundary is where that distinction is maximally informative
(crossing it is the strongest possible A $\neq$ B signal the substrate can generate).

\textbf{Consequence.}
Any substrate that can reliably detect $h \neq \mathrm{baseline}$, gate consolidation
on temporal quiescence, and propagate field state through a copy-forward step is
in the VCML class --- regardless of whether the gating mechanism is viability collapse,
a threshold potential, a chemical concentration crossing, or any other binary detector.
The lattice, the GRU, and the viability thresholds are engineering choices, not theoretical commitments.
The primitive is a difference detector with a temporal gate and a slow structural accumulator.

\subsection{VCML: the substrate}

VCML is a 2D lattice of $W \times H$ cells.
Each cell $i$ carries three state variables operating at three timescales:

\begin{center}
\begin{tabular}{llll}
\toprule
State & Symbol & Timescale & Role \\
\midrule
Hidden state   & $h_i \in \mathbb{R}^{d_h}$ & Fast (per step)          & Responds to input \\
Mid-memory     & $m_i \in \mathbb{R}^{d_h}$ & Medium (per wave)        & Accumulates contrastive signal \\
Field memory   & $F_i \in \mathbb{R}^{d_h}$ & Slow (intergenerational) & Transmitted at birth \\
\bottomrule
\end{tabular}
\end{center}

Cells have a scalar viability $v_i \in [0,1]$.
Cells with $v_i \notin [\ell, u]$ ($\ell = 0.05$, $u = 0.95$) collapse and are
immediately replaced.
At each step, spatially structured \emph{waves} perturb $v_i$: zone-class $z$
cells receive suppression (low $v$) or excitation (high $v$) depending on
their zone identity.
Standard parameters: $H=40$, $\mathrm{HALF}=40$, $N_\mathrm{zones}=4$,
zone width $z_w=10$, wave rate $\mathrm{WR}=4.8$, wave duration $= 15$ steps,
$\mathrm{FA}=0.16$, $\mathrm{FIELD\_DECAY}=0.9997$.

\subsection{VCSM: the learning rule}

Five steps, locked:

\begin{align}
\text{[Track]}       &\quad \text{baseline}_i \mathrel{+}= \beta \cdot (h_i - \text{baseline}_i) && \text{every step} \nonumber\\
\text{[Detect]}      &\quad m_i \mathrel{+}= \alpha_m \cdot (h_i - \text{baseline}_i)              && \text{during perturbation} \nonumber\\
\text{[Gate]}        &\quad m_i \mathrel{\times}= \gamma_m \quad (\gamma_m = 0.99)                 && \text{every step} \nonumber\\
\text{[Consolidate]} &\quad \text{if calm\_streak} \geq SS:\; F_i \mathrel{+}= \mathrm{FA} \cdot (m_i - F_i) && \text{on survival} \nonumber\\
\text{[Birth]}       &\quad h_\mathrm{new} = (1-\beta_s)\,\mathcal{N}(0,\sigma) + \beta_s\,F_{\text{nbr}} && \text{on collapse} \nonumber
\end{align}

\textbf{Write signal} $= h - \mathrm{baseline}$ (contrastive: change relative to recent
history).
\textbf{Gate} $=$ survival (only cells that survived the perturbation epoch consolidate).
\textbf{Copy-forward loop} $=$ the birth step: field memory propagates spatially through
inheritance from surviving neighbours.

\subsection{The copy-forward loop: formation vs.\ maintenance}

Paper~43 establishes a clean separation:

\textbf{Formation} (under sustained waves) is environmentally driven.
Zone-differentiated wave input produces differentiated viability histories, leading to
differentiated \texttt{fieldM} accumulation via any consolidation mechanism.
Even with birth seeding removed ($\beta_s = 0$) or the viability gate disabled,
sg4n remains comparable to the reference under continuous waves.
Formation is robust to all rule variants and ablations tested ($\pm50\%$ parameter changes).

\textbf{Maintenance} (after waves stop) is mechanism-specific.
The viability gate (SS parameter) is the primary load-bearing element:
without it, \texttt{fieldM} accumulates noise from unsettled cells, degrading copy-forward quality.
Removing gating drops the maintenance ratio from $1.91\times$ to $1.11\times$ (near
field-decay-only $= 0.64\times$).
Birth seeding ($\beta_s$) amplifies maintenance ($1.63\times$ without seeding) but is secondary.

\textbf{Corrected claim:} the copy-forward maintenance mechanism is the specific capability that
distinguishes VCML from simpler adaptive substrates.
It requires viability-gated consolidation as its primary structural element.

\subsection{Causal gate condition}

The quiescence gate (SS parameter) is a \emph{temporal adversarial filter}.
It is not a signal-quality filter.

This distinction is established by Papers 46--48.
The gate does not improve the amplitude of the zone signal:
a quantity-matched random write schedule ($\mathrm{rand}_{p=0.60}$, same number of writes
as SS$=10$) produces $1.69\times$ higher raw sg4 than the quiescence gate.
The gate does not improve single-cell zone specificity:
nearest-centroid zone decode accuracy (sg$_C$) is 4--5\% above chance
for \emph{all} gate conditions including SS$=0$ (Paper~48).
What the gate controls is \emph{within-zone noise}: $\sigma_w$ (mean within-zone spread
of fieldM) halves under high-SS conditions, improving the signal-to-noise ratio
$\mathrm{SNR} = \mathrm{sg4}/\sigma_w$ even as raw amplitude drops.

The gate condition and its outcome:

\begin{center}
\begin{tabular}{ll}
\toprule
\textbf{Temporal condition} & \textbf{Outcome} \\
\midrule
Site causally settled (streak $\geq$ SS) & Writes $m_i \to F_i$ \\
Site mid-perturbation (streak $< $ SS)   & Write blocked \\
Site collapses                           & New hidden state seeded from $F_i$ (copy-forward) \\
\bottomrule
\end{tabular}
\end{center}

\textbf{SS as a single performance/robustness tradeoff parameter.}
High SS delays writes until the cell has been undisturbed for $\geq$ SS steps.
This enforces two things simultaneously: (1) the write captures a settled, zone-typical
state rather than a mid-perturbation transient (reducing $\sigma_w$); and (2) the write
is blocked during adversarial input, because adversarial waves continuously restart the
calm streak counter, preventing the threshold from being reached.
Low SS permits writes during and immediately after perturbation events: higher raw amplitude,
but also higher within-zone variance, and no adversarial resistance.

The SS parameter therefore sets the operating point on a single tradeoff curve:
\begin{center}
  \textbf{high SS} $\leftrightarrow$ lower amplitude, lower noise, adversarial resistance\\
  \textbf{low SS} $\leftrightarrow$ higher amplitude, higher noise, no adversarial resistance
\end{center}
The correct causal quality metric is SNR, not raw sg4.
On SNR, the quiescence gate wins at every tested SS level (Paper~48, Exp~A).

\textbf{Gate as commitment rate controller.}
Paper~47 adds a second dimension: SS also controls the rate at which fieldM tracks
mid\_mem through the commitment sign-flip ($t^* \approx 2000$ steps, Paper~26).
High SS slows commitment (fieldM reaches the sign-flip later); low SS accelerates it.
This means the optimal $\mathrm{SS}^*$ is a function of two quantities:
adversarial pressure level and proximity to $t^*$.
Both map onto the same mechanism: the gate controls the timescale on which slow-timescale
writes couple to medium-timescale dynamics.

\textbf{Zone structure is causal, not spatial.}
A common misreading frames zone structure as a spatial phenomenon: patterns imprinted
on the lattice by spatially patterned input.
The adversarial persistence test (Paper~44) closes this reading directly.
The more precise framing is causal:
zone structure is \emph{maintained} by a causal mechanism (the copy-forward loop)
that continuously reinjects the original encoding into new cells via collapse and birth.
The zones appear to have spatial extent because the causal propagation radius
$\xi \approx r_\mathrm{wave}$ is local on the lattice --- not because the lattice
geometry is load-bearing.
Replace the lattice with a graph, a relay chain, or any other substrate with an
analogous viability-gated copy-forward step, and the mechanism is unchanged;
only the shape of the causal reach changes.
This reframes the theory's unique claim precisely:
not ``we produce spatial zone patterns'' but
``we maintain causal zone structure against adversarial environmental pressure,
via a viability gate that couples causal events (collapses) to the slow-timescale
structural manifold ($F$).''\footnote{This framing connects to the causal-set perspective
in quantum gravity (Bombelli et al.\ 1987; Sorkin 2003), where geometry emerges from
causal partial ordering rather than being fundamental.
VCML is a toy demonstration of the same direction at the level of learning dynamics:
the correlation length $\xi \approx r_\mathrm{wave}$ is a causal propagation fact,
not a geometric constant of the substrate.}

\textbf{Role of fieldM dimensionality ($d_h$): spatial vs.\ dimensional capacity.}
Memory systems can separate representational channels in two ways:
(1)~\emph{spatial separation}: assign distinct causal neighbourhoods to distinct inputs
(more cells, fixed $d_h$); or
(2)~\emph{dimensional separation}: assign distinct vector dimensions to distinct inputs
(more $d_h$, fixed cell count).
VCML uses spatial separation by design.
The HS sweep confirms that the two strategies are substitutes with diminishing returns:
HS$=5 \to$ sg4$=0.025$, HS$=2 \to$ sg4$=0.022$,
HS$=1 \to$ sg4$=0.020$ (Papers V79--V80, HS$=1$ sufficient; HS$>1$ amplifies without cliff).
The $\sim$20\% gain from HS$=1$ to HS$=5$ is the entire contribution of dimensional separation;
the remaining structure is purely spatial.

Contrast with transformer models, which use dimensional separation almost exclusively:
a 768-dimensional embedding encodes feature relationships in inner-product space,
with each ``channel'' potentially active in every context simultaneously.
The curse of dimensionality (independent causal channels conflated into Euclidean geometry)
degrades precision as feature count grows.
VCML avoids this by routing causal channels to distinct spatial locations;
the copy-forward loop ensures each location retains its causal history without
entangling it in a shared geometric space.

\textbf{Result: the irreducible primitive is a scalar thermostat.}
HS$=1$ means each cell carries a single floating-point number as its field state: a scalar
viability signal acting as a minimal difference detector.
The GRU hidden state (HS$>1$) is scaffolding that amplifies performance smoothly
via dimensional separation, but adds no discontinuous capability.
The core mechanism --- detect, gate, consolidate, transmit --- is fully operational in
the HS$=1$ limit with a single float per cell.
This is the minimum possible architecture for viability-gated spatial memory,
and it works.
The result implies that spatial separation alone is sufficient for structured memory;
dimensional expansion is a performance optimisation, not a requirement.

\textbf{Causal chain (explicit):}

\vspace{0.5em}
\begin{center}
\textit{viability perturbation (wave)}
$\;\to\;$ \textit{zone-specific viability change}
$\;\to\;$ \textit{calm-recovery produces contrastive signal}
$\;\to\;$ \textit{viability gate allows consolidation ($m \to F$)}
$\;\to\;$ \textit{collapse seeded from $F$ (copy-forward)}
$\;\to\;$ \textit{spatial propagation of zone-compatible structure}
\end{center}
\vspace{0.5em}

Each arrow is quantified by a specific VCML parameter
(wave rate, FA, SS, $\beta_s$, $\nu$).
Paper~27 confirmed the predicted scaling $t^* \sim 1/(P_c \cdot \mathrm{FA} \cdot \mathrm{WR})$
(slope $-1.0$ on log-log) before running the experiment, providing predictive validation
of this causal chain.

% ─────────────────────────────────────────────────────────────────────────────
\section{Layer I: Mechanism (Papers 0--6)}

\subsection{Death is the engine}

Mandatory turnover is not an obstacle to memory --- it is the mechanism:
\begin{equation}
\mathrm{sg4}_{C_\mathrm{ref}} > \mathrm{sg4}_{C_\mathrm{cryst}} > \mathrm{sg4}_{C_\mathrm{static}}
\end{equation}
where $C_\mathrm{ref}$ has dynamic collapses, $C_\mathrm{cryst}$ has near-zero turnover,
and $C_\mathrm{static}$ has frozen cells.
Preventing death causes crystallisation: surviving units lock into old attractors.
The optimal collapse rate is $\approx 0.004$ events per cell per step --- strictly positive.

\subsection{Three hard constraints}

\textbf{Capacity (Paper 1).} Memory density bounded by spherical random sequential adsorption:
$\rho_\mathrm{max} \propto 1/(1 + \rho r^d)$.
Geometrically hard, independent of learning rule.

\textbf{Bandwidth (Paper 2).} Effective bandwidth $\propto 1/\nu$.
Fast turnover limits signal survival to consolidation.
Tradeoff: higher $\nu$ = faster adaptation, lower bandwidth.

\textbf{Propagation (Paper 3).} Finite correlation length.
Information propagates at most one site per birth event.
Correlation length $\xi_\infty \approx 2$--3 sites (= wave radius, Paper 22).

\textbf{Noise isolation (Paper 4).} $\gamma_m < 1$ prevents false consolidation.
Without MID\_DECAY, mid-memory accumulates baseline drift and fires spuriously.
$\gamma_m$ is noise isolation, not damping.

\subsection{Memory as a formal system}

The ``memory'' terminology is operationally defined.
Three independent quantities have been measured and related to system parameters:

\begin{center}
\begin{tabular}{llll}
\toprule
\textbf{Property} & \textbf{Definition} & \textbf{Formula} & \textbf{Paper} \\
\midrule
Storage capacity  & Max distinguishable patterns & $\rho_\mathrm{max} \propto 1/(1+\rho\,r^d)$ & 1 \\
Decay rate        & sg4 half-life without input  & $\tau_{1/2} \approx 800$ steps (standard) & 18, 38 \\
Bandwidth         & Effective write rate         & $\propto P_c \cdot \mathrm{FA} / \nu$ & 2, 20 \\
\bottomrule
\end{tabular}
\end{center}

These are quantitative properties of the system, not metaphors.
Capacity is geometrically hard (spherical RSA bound, Paper~1).
Decay rate is set by FIELD\_DECAY and the metastability attractor (Paper~38).
Bandwidth is the bottleneck for how quickly new zone structure replaces old (Paper~2).

\subsection{Boundary = maximal information = maximal distinction}

The viability boundary is not a technical detail of the implementation.
It is the most informationally dense site in the substrate.

Consider what happens at a boundary crossing: cell $i$'s viability exits $[\ell, u]$.
This event is maximally informative in three senses simultaneously.

\textbf{Information-theoretic sense.}
Shannon entropy is maximised at phase transitions --- points where the probability of
being in one state versus another is close to $1/2$.
The viability boundary is precisely this: a cell near $v_i = \ell$ or $v_i = u$
is maximally uncertain between ``viable'' and ``not viable.''
The crossing event is a one-bit signal with entropy $= 1$, the maximum possible.
Interior viability states (far from the boundary) carry less information
because their future is more predictable.

\textbf{Physical sense.}
Order parameters change most rapidly at boundaries.
The gradient $|\nabla v|$ is maximal at the viability boundary.
This is where the zone-specific input signal (suppressive vs.\ excitatory wave) has had
its maximum cumulative effect on the cell's state.
The cell at the boundary has been pushed furthest from equilibrium by exactly the
perturbation history that distinguishes its zone from adjacent zones.

\textbf{VCML mechanistic sense.}
The copy-forward loop fires \emph{exclusively} at boundary crossings (collapses).
The system does not sample fieldM propagation events uniformly across the viability range;
it samples only at maximum-information events.
Memory = accumulated boundary geometry.
The fieldM of a mature VCML substrate is a compressed record of where boundary crossings
have concentrated over time --- which is precisely the record of which perturbation
patterns most strongly distinguished the zones.

\textbf{The unified statement.}
The same claim at three levels:
\begin{quote}
  (Geometry)\ \ \ \ \, Boundary $= $ maximal gradient of the order parameter.\\
  (Information) Boundary $= $ maximal Shannon entropy site.\\
  (VCML) \ \ \ \ \ \ \ Boundary $= $ maximal distinction event (A $\neq$ B at full force).
\end{quote}
These are not analogies.
They are the same fact described in the language of three different fields.
The viability boundary is the A $\neq$ B primitive at maximum expression.
The copy-forward loop is a boundary-geometry accumulator.
Zone structure is the spatial record of which boundary has been crossed most.

\subsection{Metric robustness}

The primary structure metric sg4 is the mean pairwise L2 distance between zone-mean
field vectors.
Reviewers correctly note that sg4 cannot detect polarity inversions (Papers~38, 42).
Two independent secondary metrics confirm all primary conclusions:

\textbf{nonadj/adj ratio.}
Mean L2 between non-adjacent zone pairs divided by adjacent zone pairs.
This metric carries no scale information about fieldM magnitude --- it tests only
\emph{spatial specificity} (whether farther zones differ more than nearer zones).
The nonadj/adj pattern is identical across all rule variants and ablations in Paper~43.

\textbf{mm\_lda discriminability.}
Linear discriminant analysis on mid-memory states $m_i$.
This metric reads mid-memory directly (not fieldM), providing orthogonal confirmation
that spatial information is present in the fast-timescale dynamics (Paper~30+).

All three metrics agree on all major conclusions.
The theory is not a theory of sg4; it is a theory of spatial differentiation confirmed
by three independent measurement approaches.

\subsection{ODE derivation sketch}

The zone-mean ODE (Paper~20) is derived---not fitted---from the microscopic rule:

\begin{enumerate}
\item \textbf{Averaging:} Take zone-mean of the consolidation update:
  $\Delta \bar{F} = P_c \cdot \mathrm{FA} \cdot (\bar{m} - \bar{F})$
  where $P_c = P_\mathrm{calm} \cdot P_\mathrm{consol}$ is the fraction of steps where
  consolidation occurs (product of calm-streak probability and survival probability).
\item \textbf{Decay:} Add FIELD\_DECAY: $\bar{F} \mathrel{\times}= \mathrm{FD}$.
  In continuous time: $-\nu\,\bar{F}$ where $\nu = -\ln(\mathrm{FD})$.
\item \textbf{Mid-memory:} $\bar{m}$ is driven by the wave input and damped by MID\_DECAY.
  At steady state under periodic input: $\bar{m} \approx \mathrm{FA} \cdot \delta h / (1-\gamma_m)$
  where $\delta h$ is the zone-specific deviation from baseline.
\end{enumerate}

Combining: $d\bar{F}/dt = P_c \cdot \mathrm{FA} \cdot (\bar{m} - \bar{F}) - \nu\,\bar{F}$.
Predicted growth rate agrees with simulation within 3\% (Paper~20).
The spatial PDE (Papers~22--26) extends this by retaining the spatial $\nabla^2$ and $(\nabla F)^2$
terms that average to zero in the zone-mean ODE.

\subsection{AI failure modes (Paper 6)}

Every major AI memory architecture maps to a VCML failure mode:

\begin{center}
\begin{tabular}{lll}
\toprule
Architecture & Failure mode & VCML equivalent \\
\midrule
Transformer (context window) & Evaporation: no consolidation     & $\nu \to \infty$ \\
RAG                          & Bandwidth: retrieval misses        & $P_\mathrm{calm}$ too low \\
Continual learning           & Catastrophic forgetting            & No $\gamma_m$ (false positives) \\
RNN                          & Crystallisation: fails to update   & $\nu \to 0$ \\
\bottomrule
\end{tabular}
\end{center}

These failures are \emph{structural} constraints, not engineering defects.
The three unavoidable failure modes are: \textbf{bias} (consolidation is exposure-weighted),
\textbf{confabulation} (recall reconstructs from $F$ fragments),
and \textbf{forgetting} ($\mathrm{FIELD\_DECAY} < 1$ is required for adaptivity).

\textbf{The ontological stack and the positive claim.}
VCML is not merely a model of hippocampal memory.
It is an instantiation of the most minimal computing primitive: the \emph{difference
detector} (thermostat) --- a device that registers ``perturbed vs.\ calm'' and
adjusts a persistent variable accordingly.
Every more powerful computing paradigm is the thermostat with additional properties enabled:

\begin{center}
\begin{tabular}{ll}
\toprule
Layer & What is added over the layer below \\
\midrule
Difference detector (thermostat) & Persistent variable; responds to a single signal \\
Memory                           & Multiple slots; time-indexed write/read \\
Relations                        & Associations between slots (structured $F$) \\
Dynamics                         & Time-varying relational structure (sequences) \\
Symbolic reasoning               & Stable attractors that act as discrete tokens \\
\bottomrule
\end{tabular}
\end{center}

The failure-mode table above is the complementary view: each AI architecture
ablates one or more layers of this stack.
The VCML adaptive regime occupies the interior --- not maximum memory (crystallised),
not maximum plasticity (sub-viability), but the driven steady state in between.

\textit{Note:} This stack is a post-hoc organisational framework, not a derived result.
The claim that ``every computing paradigm is the thermostat with properties ablated'' is
a motivated framing that organises what Papers~0--43 found; it is not itself proven by
those papers.
The formal criterion for what counts as ``the thermostat with property $X$ ablated'' ---
as opposed to a different primitive entirely --- remains an open question.
The failure-mode table is a mapping, not a derivation.

% ─────────────────────────────────────────────────────────────────────────────
\section{Layer II: Parameter Geometry (Papers 7--21)}

\subsection{Three behavioural regimes}

\begin{center}
\begin{tabular}{lll}
\toprule
Regime & Condition & Signature \\
\midrule
Sub-viability & $\nu$ too high & Cells die before consolidating. $F \approx 0$. sg4 $\to 0$. \\
Adaptive & $\nu$ in window & Dynamic spatial structure via copy-forward. \\
Crystallised & $\nu$ too low & $F$ frozen. Cannot update. \\
\bottomrule
\end{tabular}
\end{center}

\subsection{Analytic adaptive window}

\begin{align}
\nu_\mathrm{cryst} &\approx \mathrm{FD}^{1/\nu_\mathrm{cryst}} && \text{field decay kills memory before update} \\
\nu_\mathrm{max}   &\approx P_\mathrm{calm} \cdot \mathrm{FA} \cdot \mathrm{FD}^{1/\nu} && \text{signal must survive at least one lifetime} \\
\nu^*              &= \sqrt{\nu_\mathrm{cryst} \cdot \nu_\mathrm{max}} && \text{geometric mean (empirically, factor 2)}
\end{align}

\subsection{Amplitude laws}

\textbf{Saturation law (Paper 16):}
\begin{equation}
\mathrm{sg4}(\mathrm{FA}) \approx \frac{C \cdot \mathrm{FA}}{\mathrm{FA} + K_\mathrm{eff}}, \qquad
K_\mathrm{eff} = \frac{\kappa}{P_\mathrm{consol}} \approx 0.114
\end{equation}
The apparent power-law exponent 0.43 (Paper~15) is the local slope of this saturation curve.

\textbf{Two-factor decomposition (Paper 17):}
\begin{equation}
\mathrm{sg4} = G \cdot R, \quad
G = \Gamma \cdot K_\mathrm{eff}(\text{zone\_w}), \quad
R = \frac{\mathrm{FA}}{\mathrm{FA} + K_\mathrm{eff}}
\end{equation}

\textbf{Amplitude non-monotonicity (Papers 15, 32):}
sg4 peaks at collapse rate coll/site $\approx 0.004$.
Over-perturbation (coll/site $> 0.01$) destroys structure.
Optimal amplitude threshold $A^* \approx 0.30$ is endogenous (set by internal noise floor).

\subsection{ODE closure}

Zone-mean field obeys (Paper~20):
\begin{equation}
\frac{d\bar{F}}{dt} = P_c \cdot \mathrm{FA} \cdot (\bar{m} - \bar{F}) - \nu\,\bar{F},
\qquad P_c = P_\mathrm{calm} \cdot P_\mathrm{consol}
\end{equation}
Predicted growth rate within 3\% of simulation. A spatial formation delay
(not in ODE) arises from birth-mediated propagation (Papers 20--21).

\subsection{Scaling (V92--V94)}

Dynamic advantage (+43\% to +70\%) holds at all tested scales ($N = 1{,}600$ to $102{,}400$).
Optimal WR $\sim 2 \times N_\mathrm{active}$ (Poisson scaling).
sg4 raw declines with scale because of finite correlation length $\xi \approx$ wave radius.
Beyond $\xi$ the copy-forward loop cannot maintain zone-wide coherence.

\textbf{Self-organization result: the substrate spontaneously partitions into regional patches.}
At large scale ($N \geq 25{,}600$), the finite correlation length has an important consequence:
rather than failing to form structure, the substrate forms \emph{coherent local patches} from
a blank grid initialization, with each patch maintaining its own zone-specific fieldM
independently of patches beyond distance $\xi$.
This is not a failure of the mechanism --- it is the mechanism operating at its natural spatial scale.
The resulting multi-patch landscape is not externally imposed; it self-organizes from
locally identical rules, wave input, and mandatory turnover.
At the right spatial scale (zone width $\sim \xi$) the system forms a regional architecture
spontaneously: a blank grid, a structured wave, and the copy-forward loop are sufficient.

% ─────────────────────────────────────────────────────────────────────────────
\section{Layer III: Spatial PDE (Papers 22--26)}

\subsection{Field equation}

The field $F$ obeys a stochastic hybrid PDE:
\begin{equation}
\frac{\partial F}{\partial t} = P_c \cdot \mathrm{FA} \cdot (\bar{m} - F)
+ \kappa \nabla^2 F + \eta_\mathrm{copy} + \eta_\mathrm{birth}
\end{equation}

Two universality class regimes determined by $\nu/\kappa$:

\begin{center}
\begin{tabular}{lll}
\toprule
Regime & Governing equation & Condition \\
\midrule
Allen-Cahn & $\partial F/\partial t = \kappa\nabla^2 F - \lambda(F-F^*)^2 + \eta$ & $\nu/\kappa \ll 1$ \\
Burgers    & $\partial F/\partial t = \kappa\nabla^2 F + \mu(\nabla F)^2 + \eta$ & $\nu/\kappa \gg 1$ \\
\bottomrule
\end{tabular}
\end{center}

The Burgers term arises from the copy-forward birth step introducing a $\nabla F$ nonlinearity.
At the transition ($\nu/\kappa \approx 1$), the effective advection term
\emph{reverses sign} (Burgers reversal, Papers 24--26): copy-forward sharpens zone
boundaries from the outside rather than from the inside as Allen-Cahn would.

\subsection{Interface theory}

Zone interfaces have width $\approx 2$ sites = wave radius $r_w$, not set by $\kappa$
or interface tension.
They undergo a biased random walk driven by wave input, not gradient descent (Paper~28).

\subsection{Optimal zone width}

\begin{equation}
W_\mathrm{zone}^* \sim \sqrt{\kappa/\nu} \quad (\approx 4.5 \text{ at standard params})
\end{equation}
Balances diffusion smearing (wider zones benefit from $\kappa$ smoothing) against
turnover noise (Paper~26).

% ─────────────────────────────────────────────────────────────────────────────
\section{Layer IV: Dynamics (Papers 27--33)}

\subsection{Commitment epoch}

Zone structure does not form monotonically.
The system first adopts the \emph{wrong} spatial polarity (anti-correlated with final state),
then flips through a commitment epoch:
\begin{equation}
t^* \approx 1600\text{--}2000 \text{ steps (standard params)}, \qquad
t^* \sim \frac{1}{P_c \cdot \mathrm{FA} \cdot \mathrm{WR}}
\end{equation}
Mechanism is mixed pitchfork--spinodal: zones tend to flip together within a seed
(pitchfork), but the flip time varies across seeds (spinodal component).

\subsection{Attractor structure (Paper 32)}

The zone attractor is a diffuse noisy cloud (CV $\approx 0.44$), not a fixed point.
The system is bimodal: $P_\mathrm{high}(\mathrm{geo}) = 0.50$ (geographic coupling) vs.\
$P_\mathrm{high}(\mathrm{ctrl}) = 0.27$ (no coupling).
Basin depth $G_\mathrm{basin} = 1.88$: the high-structure basin has a substantial
free-energy well.
No Kramers escape events observed (no sudden zone swaps at steady state).

\subsection{Formal non-gradient proof (Paper 33)}

\textbf{Theorem} (machine-verified in Lean 4 + Mathlib4, 3103 jobs, 0 errors):\\[0.3em]
VCML is \emph{not} a conservative Euclidean gradient system.

\begin{equation}
\text{curl}(f_F, f_m) = \frac{\partial f_m}{\partial F} - \frac{\partial f_F}{\partial m}
= 0 - \mathrm{FA} = -\mathrm{FA} \neq 0
\end{equation}

where $f_F = \mathrm{FA}(m - F)$ and $f_m = -(1-\gamma_m)m$.
By Clairaut's theorem, no $C^2$ potential $V$ exists with $\nabla V = -f$.
\emph{Honest caveat:} rules out Euclidean gradient flow only; mirror descent,
Riemannian, and nonlocal potentials are not ruled out.

\textbf{Physical interpretation:} consolidation is asymmetric ($F$ chases $m$, $m$ does
not chase $F$); energy is injected via waves and dissipated via field decay.
The system is a fundamentally driven non-equilibrium process.

% ─────────────────────────────────────────────────────────────────────────────
\section{Layer V: Multi-Region Coupling (Papers 34--39)}

\subsection{Phase diagram (Paper 34)}

The Allen-Cahn $\to$ Burgers transition has measurable critical exponents.
The two manifolds $C_1$ (viability volume) and $C_2$ (structural manifold) share $\beta_s$,
but are otherwise independent in the measured parameter ranges.

\subsection{Geographic relay coupling (Papers 35--37)}

Two separated regions $A$, $B$ can communicate zone structure via geographic relay:
region~$B$ receives the same wave sources as region~$A$, producing correlated
(not identical) fieldM patterns.

\textbf{Relay gain:}
$G = \mathrm{sg4}(B) / \mathrm{sg4}(A)$, measured at matched parameters.

\textbf{Two conditions for relay to succeed:}
\begin{align}
N < N_\mathrm{crit} &= \frac{L}{5\,r_w} && \text{(resolution: relay grid not too fine)} \\
\Phi &= \frac{\mathrm{WR} \cdot r_w}{T_\mathrm{dur}} < 1 && \text{(coverage: waves don't saturate)}
\end{align}

\textbf{Correlation length:} $\xi_\mathrm{eff} = r_w = 2$ sites (wave-dominated, independent
of $\kappa$, Paper~36). The zone map is replicated at resolution $\approx r_w$, not $\xi_\mathrm{diffusion}$.

\textbf{Chain relay (Paper 37):} relay gain $G_k \approx G_1$ is depth-stable
(no accumulation or collapse along a relay chain of arbitrary depth).
Empirical coefficient: $N_\mathrm{crit} = L/(5\,r_w)$ (coefficient 5, not 4).

\textbf{$r_w$ scaling:} $N_\mathrm{crit} \propto 1/r_w$ confirmed across $r_w \in \{1,2,3\}$.
At $r_w = 4$: $\Phi > 1$ (over-perturbation), $G < 1$ everywhere.

\subsection{Cleft geometry and alignment blindness (Paper 39)}

When a synaptic cleft (spatial gap) exists between source and relay regions,
\emph{deterministic} offsets in zone labelling do not degrade relay gain $G$
regardless of magnitude (at $D_\mathrm{cleft} = 100\%$ misalignment, $G = 3.42$).
\emph{Stochastic} noise does degrade $G$ by randomising zone identity across waves.
$G_\mathrm{det}/G_\mathrm{stoch} = 1.25$ on average at matched magnitudes.

\textbf{Principle:} \emph{consistency} of geographic routing is the operative
requirement, not geometric accuracy.
Systematic errors are tolerated; random errors are not.

\textbf{Biological prediction:} topographic inversion (consistent shift) produces
normal downstream zone formation; scattered terminal fields (high noise) impair it.

% ─────────────────────────────────────────────────────────────────────────────
\section{Layer VI: Metastability, Continual Learning, and Robustness (Papers 38--43)}

\subsection{Metastability (Paper 38)}

Under continuous single-task input:
\begin{itemize}[itemsep=2pt]
  \item Zone structure commits at $t^* \approx 2000$ steps.
  \item After commitment, sg4 decays with \textbf{half-life $\approx 800$ steps}.
  \item By $T = 6000$--$8000$ steps, sg4n falls below 0.10.
\end{itemize}
The bimodal attractor has $P_\mathrm{high}(\mathrm{geo}) = 0.50$, $G_\mathrm{basin} = 1.88$.
Polarity inversion mid-run is invisible to sg4 (polarity-blind metric).

\subsection{Real continual learning (Paper 40)}

\textbf{Permutation transparency:} a cyclic zone shift ($\Delta = z_w$) is immediately
transparent --- sg4n$_\mathrm{new}$ = sg4n$_\mathrm{old}$ at the shift moment.
sg4 is permutation-invariant.
Genuine continual learning requires spatial reassignment \emph{with overlap}.

\textbf{Partial-overlap disruption:} a half-width shift ($\Delta = z_w/2$) drops sg4n$_\mathrm{new}$
by 50\% at the shift, then recovers slowly ($> 2500$ steps to recover; still rising at $T=4000$
where it reaches sg4n $= 0.80$, exceeding the C$_\mathrm{ref}$ peak of $0.38$).

\textbf{Gradual forgetting:} post-commitment interleaving (Exp~C) drives both tasks to equal
steady-state sg4n $\approx 0.28$--$0.29$, regardless of prior commitment.
Forgetting is gradual (rate set by metastability half-life $\approx 800$ steps),
not catastrophic.

\subsection{Copy-forward refreshing and task switching (Paper 41)}

\textbf{Main result:} the predicted threshold $T_\mathrm{crit} \approx 800$ steps does not exist.
Every tested switching period maintains sg4n $5$--$10\times$ above the single-task baseline
at $T = 6000$ (C$_\mathrm{ref}$ sg4n $= 0.063$; all interleaved sg4n $= 0.30$--$0.64$).

\begin{center}
\begin{tabular}{rccc}
\toprule
$T_\mathrm{switch}$ & sg4n$_A$ & sg4n$_B$ & Amplification \\
\midrule
50   & 0.563 & 0.288 & $8.9\times$ \\
500  & 0.615 & 0.546 & $9.7\times$ (optimum) \\
800  & 0.639 & 0.338 & $10.1\times$ \\
3000 & 0.456 & 0.507 & $7.2\times$ \\
C$_\mathrm{ref}$ & 0.063 & --- & $1.0\times$ \\
\bottomrule
\end{tabular}
\end{center}

\textbf{Mechanism --- copy-forward refreshing:}
task switches trigger zone-boundary collapses $\to$ births $\to$ fieldM replenishment.
This continuously counteracts FIELD\_DECAY. The metastability half-life (800 steps) is a
\emph{static} decay constant applying only when no perturbation is present.
Any periodic perturbation resets the decay clock.

\textbf{Optimal switching period:}
$T_\mathrm{switch}^* \approx 0.6 \times \tau_{1/2} \approx 500$ steps.
Non-monotone: too short (no commitment within epoch); too long (decay between switches).

\textbf{Biological implication:}
sleep replay (``task switching'' during consolidation) maintains hippocampal zone structure
via mandatory turnover.
Absence of replay $\to$ structural decay on the metastability timescale ($\sim 800$ steps).
Not just encoding --- structural maintenance.

\subsection{Three anomalies closed (Paper 42)}

\textbf{Anomaly 1 (C$_\mathrm{loc}$ ratio $> 1$, V89).}
Waves are not necessary for copy-forward maintenance.
When waves stop at the commitment peak ($T = 2000$), structure decays with ratio $0.664$
vs.\ field-decay-only prediction of $0.549$ --- structure persists \emph{above} pure field
decay because ageing and instability noise generate $\approx 2.84$ collapses/step (spontaneous
copy-forward).

\textbf{Anomaly 2 (C$_\mathrm{half}$ late surge).}
The surge in Paper~40 (delta-5 shift reaching sg4n $= 0.80$) is sustained.
At $T = 6000$, C$_\mathrm{half}$ maintains sg4n $= 0.212$ vs.\ C$_\mathrm{ref} = 0.068$
(3.1$\times$).
Mechanism: zone-boundary mismatch drives sustained structured collapses $\to$ continuous
copy-forward refresh.
Identical in kind to Paper~41's task-switching amplification.

\textbf{Anomaly 3 (T$_\mathrm{sw}=50$/100 apparent reversal).}
The reversal (B$>$A at $T_\mathrm{sw}=100$) is a \emph{checkpoint phase aliasing artifact}.
$300/T_\mathrm{sw}=3$ (odd): the last three checkpoints fall 2:1 in task~B epochs,
inflating sg4n$_B$.
Epoch-end measurement (unbiased) restores sg4n$_A \geq$ sg4n$_B$ for all switching periods.
\textbf{Methodological lesson:} use epoch-end measurements in all switching experiments.

\subsection{Rule-family robustness (Paper 43)}

Paper~43 directly addresses the reviewer concern: ``is the phenomenon robust to nearby
rule changes, or is it fragile rule engineering?''

\textbf{Experiment A} (continuous waves, 10 conditions $\times$ 5 seeds):
All conditions---including ablations of birth seeding and viability gating---produce
sg4n comparable to the reference.
Formation is environmentally driven, not rule-specific.

\textbf{Experiment B} (waves stop at $T{=}1500$, 4 conditions $\times$ 5 seeds):

\begin{center}
\begin{tabular}{lrrl}
\toprule
\textbf{Condition} & \textbf{Maintenance ratio} & \textbf{Field-decay-only} & \textbf{Interpretation}\\
\midrule
\texttt{ref} (all mechanisms)    & 1.91 & 0.64 & strong maintenance\\
\texttt{no\_birth\_seed} (SB=0) & 1.63 & 0.64 & amplifier removed\\
\texttt{no\_gating} (SS removed) & 1.11 & 0.64 & near field-decay-only\\
\texttt{no\_both}                & 1.65 & 0.64 & floor set by gating\\
\bottomrule
\end{tabular}
\end{center}

Viability gating is the primary load-bearing element.
Birth seeding is a secondary amplifier.
All six rule variants ($\pm50\%$ perturbations of SS, MID\_DECAY, SEED\_BETA) maintain
comparable sg4n to the reference, confirming robustness.

\textbf{Complete rule ablation table:}

\begin{center}
\begin{tabular}{lll}
\toprule
\textbf{Component removed} & \textbf{Formation} & \textbf{Maintenance} \\
\midrule
Waves (WR$=0$, Paper 42)          & structure persists (scale-invariant) & spontaneous copy-forward persists \\
Birth seeding ($\beta_s{=}0$, P43) & robust ($\sim$ref sg4n)             & reduced ($1.63\times$ vs $1.91\times$) \\
Viability gating (SS$=0$, P43)    & slightly reduced ($0.20$ vs $0.27$) & strongly reduced ($1.11\times$) \\
Both mechanisms (P43)             & robust                               & $1.65\times$ (floored by gating) \\
fieldM altogether (Papers 0--5)   & collapses to zero                   & not applicable \\
\midrule
\multicolumn{3}{l}{\textit{Rule variants: all maintain sg4n $\geq 50\%$ of ref under waves}}\\
SS perturbed $\pm50\%$        & robust & --- \\
MID\_DECAY perturbed $\pm50\%$ & robust & --- \\
SEED\_BETA perturbed $\pm50\%$ & robust & --- \\
\bottomrule
\end{tabular}
\end{center}

\subsection{Adversarial persistence: memory vs.\ imprinting (Paper 44)}

Paper~44 addresses the hostile-reviewer objection directly:
\emph{``The system merely mirrors environmental structure; what is missing are experiments
demonstrating structure maintenance under conditions where simple imprinting would fail.''}

\textbf{Design:} All three conditions (adaptive, passive, rigid) use the \emph{identical}
full copy-forward mechanism during encoding ($T = 0$--$3000$, past the commitment epoch).
At $T_\mathrm{encode}$, wave-zone assignments are \emph{permanently flipped} (even zones
receive excitation, odd zones receive suppression --- opposite of encoding).
Adversarial amplitude is 50\% of encoding amplitude to allow competition to unfold gradually.
The passive condition switches off viability gating and birth seeding;
the rigid condition removes cell turnover entirely.

\textbf{Primary metric:} cosine fidelity $= \frac{1}{4}\sum_{z=0}^{3}
\cos(\bar{F}_z(t),\,\bar{F}_z(T_\mathrm{encode}))$.
Positive = original structure preserved; negative = adversarial pattern encoded.

\begin{center}
\begin{tabular}{lrrrl}
\toprule
Condition & Fid@200 & Fid@800 & Fid@2000 & Interpretation \\
\midrule
Adaptive  & $-0.111$ & $-0.089$ & $+0.129$ & Memory: resists, recovers \\
Passive   & $-0.272$ & $+0.006$ & $-0.076$ & Imprinting: encodes adversarial \\
Rigid     & $-0.155$ & $-0.452$ & $+0.167$ & Frozen, then decay-to-prior \\
\bottomrule
\end{tabular}
\end{center}

\textbf{Key results:}
\begin{enumerate}[itemsep=1pt]
  \item Late-phase fidelity ($T_\mathrm{adv}=2000$): adaptive $+0.129$,
        passive $-0.076$, $\Delta = 0.205$,
        $t \approx 8.9$, $p < 0.001$.
        A pure imprinting system converges to adversarial statistics under sustained
        adversarial input; the adaptive mechanism does not.
  \item Removing turnover (rigid condition) \emph{increases} intermediate-timescale
        adversarial adaptation (fidelity $-0.452$ at $T=800$ vs adaptive $-0.089$).
        The copy-forward loop is the resistance mechanism; disabling it removes resistance.
  \item Adaptive fidelity recovers to positive at long timescales: the original encoding
        is a dynamical attractor maintained by the copy-forward loop, not a static trace.
\end{enumerate}

Papers~43 and~44 jointly close the necessity--and--resistance question:
Paper~43 shows viability-gating is necessary (ablation test);
Paper~44 shows the mechanism actively resists adversarial input
(the test that distinguishes memory from imprinting).

\subsection{External probe: causal geometry in GPT-2 (Paper 45)}

Paper~45 tests the causal geometry conjecture outside VCML using GPT-2 as a probe system.
If representation geometry is organised by causal coupling (not Euclidean proximity),
then tokens that causally influence each other should have closer representations
in the layers where causal integration occurs---after controlling for sequence distance.

\textbf{Method.}
Mean-ablation perturbation (replacing token $i$'s embedding with the vocabulary mean)
measures ablation influence $A_{ij}^\ell$ on token $j$ at layer $\ell$.
Geometric distance is $G_{ij}^\ell = 1 - \cos(\mathbf{h}_i^\ell, \mathbf{h}_j^\ell)$.
Partial correlation $r_p^\ell = \mathrm{pcorr}(G, A \mid d_{ij})$ controls for sequence
distance on 11{,}706 token pairs from 50 UD English EWT sentences.

\textbf{Results.}
\begin{enumerate}
  \item \textbf{Sign flip at layer~4.}
    Layers~1--3: positive $r_p$ (more influence $\to$ more distant, feature construction).
    Layer~4: sign flip to negative.
    Layers~5--11: negative $r_p$, peaking at layer~7
    ($r_p = -0.337$, $p \approx 10^{-308}$, feature integration).
    Layer~12: reverts to positive (output re-expansion).

  \item \textbf{Random-token control tracks ablation.}
    Replacing token $i$ with a fixed random token yields $r_p = -0.331$ at layer~7---
    nearly identical to the ablation profile.
    The signal is in the attention routing structure, not in the content of the substitution.
    This is the transformer analog of Paper~36's $\xi \approx r_w$ invariance: the
    geometric effect is in the causal routing, not in the transported content.

  \item \textbf{Dependency-distance matched.}
    Syntactically linked pairs are geometrically closer than unlinked pairs at matched
    sequence-distance bands ($p < 0.05$ in all three bands: $[1,3)$, $[3,6)$, $[6,15)$).
    The effect is largest at the $[3,6)$ band where sequence-distance confounds are weakest.

  \item \textbf{Mixed-effects regression.}
    Ablation influence is a significant negative predictor of geometric distance at layer~6
    ($\hat{\beta} = -0.00309$, $p \approx 10^{-196}$).
    Syntactic dependency label is \emph{not} an independent predictor beyond ablation
    influence ($p = 0.509$): syntax is mediated by causal routing, not separately encoded.
\end{enumerate}

\textbf{Implication for VCML.}
GPT-2's representational geometry is organised by causal coupling in exactly the regime
(middle layers) where integration replaces feature construction.
This supports the causal reparametrisation of $C_2$: correlation length is set by causal
reachability, not diffusion length.
The spatial-vs.-dimensional capacity distinction is confirmed: VCML's spatial separation
at HS$=1$ avoids the channel-conflation problem inherent in transformer-style dimensional
separation ($d_h = 768$).

\subsection{Temporal gate and $\Omega$ validation (Paper 46)}

Paper~46 closes two claims arising from the causal geometry programme.

\textbf{Exp~A (temporal gate): the gate is an adversarial filter.}
Six conditions varying the consolidation gate --- SS$\in\{0,5,10,15,20\}$ and a
random-write control (rand\_p60, $p=0.60$) at matched rate --- reveal an unexpected pattern:
random writes outperform quiescence-gated writes by $1.68\times$ in benign conditions
(rand\_p60 sg4$=0.041$ vs.\ SS$=10$ sg4$=0.024$, $t=-4.32$, $p=0.003$).
All quiescence-gated conditions cluster tightly ($0.022$--$0.025$) despite varying
write rates; both ungated conditions reach $0.041$.

The mechanism: mid\_mem accumulates zone-specific signal during and immediately after
wave perturbation.
The quiescence gate waits SS$=10$ steps before writing, by which time mid\_mem has
partially decayed ($0.99^{10} \approx 0.904$) and the peak-signal writes have been excluded.
Random writes at $p=0.60$ capture a fraction of the peak-signal perturbation phase, boosting sg4.

Reconciliation with Paper~44: under adversarial waves, mid\_mem during perturbation encodes
the wrong-zone signal; waiting SS steps allows this signal to decay before writing,
protecting the original encoding.
The gate is an \emph{adversarial filter}: it reduces steady-state structural strength
($1.68\times$) but is essential for adversarial robustness.
The optimal SS threshold is set by adversarial pressure, not by signal quality.

\textbf{Exp~B ($\phi_w$ validation): $\Omega = \mathrm{WR}\times\phi_w$ is the operative variable.}
Three wave radii ($r_w \in \{1,2,3\}$, $\phi_w \in \{5,13,25\}$) under two WR regimes:
$\Omega$-constant (WR adjusted so $\mathrm{WR}\times\phi_w=62.4$) and WR-constant (4.8).
The $\Omega$-constant sweep is approximately flat (CV$=0.126$, $<0.15$ threshold)
for $r_w \in \{2,3\}$; modest deficit at $r_w=1$ attributed to $\delta = W_\mathrm{zone}/r_w = 10$
(zone too wide relative to causal horizon for single-event propagation to maintain coherence).
The WR-constant sweep increases monotonically with $r_w$ ($0.012 \to 0.024 \to 0.031$),
confirming that neither WR nor $r_w$ alone is operative.
First direct experimental support for the $C_2$ causal reparametrisation.

\subsection{$\delta$-sweep and gate--commitment interaction (Paper 47)}

Paper~47 completes the $C_2$ causal reparametrisation validation.

\textbf{Exp~A ($\delta$-sweep): third $C_2$ coordinate confirmed.}
Five zone configurations sweep $\delta = W_\mathrm{zone}/r_\mathrm{wave}$
across $\{10, 5, 4, 2.5, 2\}$ at fixed $\Omega = 62.4$.
The raw sg4 increases with more zones (small $\delta$), partly a counting
artefact (more zone pairs).  The encoding-quality signal is the
non-adjacent/adjacent ratio $R_\mathrm{na}$:
\begin{itemize}
  \item $\delta = 5$ (standard 4-zone grid): $R_\mathrm{na} = 1.061$
        (monotonic zone gradient, nonadj $>$ adj).
  \item $\delta = 4$--$2.5$: $R_\mathrm{na} \approx 0.97 < 1$
        (adjacent zones more distinct than non-adjacent; wave-level
        rather than zone-level encoding).
  \item $\delta = 2$: $R_\mathrm{na} = 1.04$, slight recovery from even/odd
        polarity reinstating a two-class structure at sub-zone grain.
\end{itemize}
The viable range is $\delta \geq 5$: below this, the wave footprint spans
zone boundaries and the copy-forward loop builds sub-zone spatial
gradients rather than monotonic zone-level ones.
The fraction of home-zone events for a boundary cell scales as
$\delta/(\delta+1)$; at $\delta=5$ this is $83\%$, falling to $67\%$ at
$\delta=2$ (zone discrimination threshold crossed).

\textbf{Exp~B (gate~$\times$~adversarial): commitment rate, not resistance.}
A two-phase protocol with $T_\mathrm{encode}=2000$ baseline snapshot
crosses gate conditions ($\mathrm{SS}\in\{0,5,10,15,20,\mathrm{rand}_{0.60}\}$)
with adversarial amplitude ($a\in\{0,0.125,0.25,0.50\}$).
A striking result: at $a=0$ (no adversarial), fidelity relative to the
$T=2000$ snapshot is near-zero or negative for most conditions.
SS$=0$ reaches $-0.22$ under normal (non-adversarial) waves.

This reflects the commitment epoch (Paper~26): fieldM at $T=2000$
is pre-committed for many seeds, encoding an anti-correlated
pre-pattern.  Low-SS gates accelerate commitment --- fieldM crosses the
sign flip and ends anti-correlated with the baseline.
High-SS gates slow commitment --- fieldM remains near the $T=2000$
state, yielding positive fidelity.

\textbf{Synthesis.}  Papers 26, 46, and 47 describe one mechanism:
the quiescence gate controls the timescale on which fieldM tracks
mid\_mem.  At pre-commitment baselines, this manifests as commitment
rate.  Under adversarial input, this manifests as adversarial filtering.
The optimal SS$^*$ is a function of both adversarial pressure level
and proximity to $t^*$: a high-pressure/post-commitment environment
calls for high SS; a low-pressure/pre-commitment environment
benefits from lower SS to complete the sign flip on schedule.

\subsection{Population codes and the causal structure metric $\mathrm{sg}_C$ (Paper 48)}

Paper~48 resolves a methodological inconsistency: the primary metric sg4 is
population-level (mean pairwise L2 between zone means), but we had not established
whether zone information lives at the single-cell or population level.

\textbf{Population code finding.}
Single-cell zone decoding accuracy (sg$_C$, nearest-centroid classification in
fieldM space) is only 4--5\% above chance across all gate conditions
($N_\mathrm{zones}=4$, chance$=0.25$).
Zone structure is not carried by individual cells --- it is distributed across
the population, recoverable from zone means but not from any individual fieldM vector.
This is architecturally expected: mandatory turnover erases individual-cell labels
by construction.  Population-level labels persist because they are carried by the
mean, not by any particular unit.

\textbf{SNR as causal quality measure.}
Define $\sigma_w$ as the mean within-zone standard deviation of fieldM, and
$\mathrm{SNR} = \mathrm{sg4}/\sigma_w$.
The quiescence gate's operative effect is \emph{noise reduction}: at high SS,
$\sigma_w$ halves while sg4 drops by less.  The SNR ordering is:
$\mathrm{SS}{=}20\ (0.129) > \mathrm{SS}{=}10\ (0.125) > \mathrm{rand}_{0.60}\ (0.108) > \mathrm{SS}{=}0\ (0.105)$.
This is the correct causal ordering.  The Paper~46 finding that
$\mathrm{rand}_{0.60}$ achieves 1.69$\times$ higher sg4 than SS$=10$ is now
explained: the random gate has high amplitude but $2\times$ the within-zone noise;
the quiescence gate trades amplitude for precision.

\textbf{Direction validation.}
In the $\delta$-sweep, sg$_C$ decreases monotonically as $\delta$ shrinks
($0.067 \to 0.022$, correct causal direction), while sg4 increases
(counting artefact from more zone pairs, wrong direction).
sg$_C$ and na\_ratio agree: $\delta \geq 5$ is the threshold.

\textbf{Full metric suite.}  Future experiments should report:
sg4 (between-zone signal), $\sigma_w$ (within-zone noise),
SNR (causal quality), na\_ratio (encoding geometry),
sg$_C$ (direction check when $N_\mathrm{zones}$ varies).

\subsection{$\beta_s$, gate optimum, and population-code relay (Paper 49)}

Paper~49 closes three open threads from Papers~46--48.

\textbf{$\beta_s$ is a maintenance coordinate, not a formation coordinate.}
A full sweep of $\beta_s \in \{0, 0.1, 0.25, 0.5, 0.75, 1.0\}$ at fixed
$\Omega=62.4$ and $\delta=5$ produces flat sg4 ($0.0195$--$0.0202$) and flat
SNR ($0.100$--$0.104$) --- all within noise.
Under continuous wave input, the wave-driven viability differentiation overwrites
the birth seed state within $\sim$10--50 steps, masking any $\beta_s$ effect.
The C$_2$ triplet is complete: $\Omega$ and $\delta$ are formation coordinates;
$\beta_s$ is a maintenance coordinate (active only when waves are absent,
consistent with Paper~43's $1.63\times$ maintenance degradation at $\beta_s=0$).

\textbf{Optimal gate reverses at the commitment epoch.}
The 2D surface SS$^*(T_\mathrm{enc}, a)$ has a clean structure:
before $t^*$ ($T_\mathrm{enc}=1000$), SS$=0$ wins under adversarial pressure
--- rapid commitment is more valuable than adversarial resistance pre-commitment;
after $t^*$ ($T_\mathrm{enc}\geq 2000$) with strong adversarial ($a=0.5$), SS$=10$
wins --- the quiescence gate's adversarial resistance dominates;
after $t^*$ without adversarial, SS$=20$ wins --- noise reduction is the objective.
Weak adversarial ($a=0.25$) favours SS$=0$ throughout --- the gate is too
conservative for low-pressure environments.
The gate reversal at $t^*$ unifies Papers~26, 46, 47, and~49:
optimal $\mathrm{SS}^* = f(T_\mathrm{enc}/t^*,\, a)$, not a constant.
A new dissociation also emerges: under strong post-commitment adversarial input,
SNR and decode\_norm (sg$_C$, per-cell location accuracy) diverge.
SS$=0$ achieves higher sg$_C$ (per-cell accuracy preserved by continuous rewrites
that track the adversarial signal) while SS$=10$ achieves higher SNR (zone-mean
separation protected by blocking adversarial rewrites).
Population-level geometry and individual-cell accuracy are distinct objectives
with opposite optimal gate strategies under adversarial pressure.

\textbf{Relay requires structured ongoing viability events, not only structured seeding.}
Experiment~C compares three L2 coupling modes:
mean\_relay (births seeded from L1 zone-mean fieldM, random waves to L2),
geo (same wave amplitude pattern as L1, Paper~35 relay),
ctrl (random waves, own fieldM).
Result (with full metric suite including sg$_C$/decode\_norm and na\_ratio):
ctrl ($1.00\times$, na\_ratio$=1.245$) $>$ geo ($0.89\times$, na\_ratio$=0.826$)
$>$ mean\_relay ($0.80\times$, na\_ratio$=0.908$).
\emph{Neither} relay mode outperforms the baseline.

\textit{Why geo fails.}
The wave-class alternating pattern (even zones suppressed, odd zones excited)
creates zone \emph{parity} encoding rather than zone \emph{identity} encoding.
Adjacent zones receive opposite perturbations; non-adjacent zones receive similar ones.
The na\_ratio $<1$ signature directly reveals this: adjacent zone means are
more separated than non-adjacent zone means.

\textit{Why mean\_relay fails.}
L2's ongoing fieldM updates ($\Delta$fieldM $\propto$ mid$-$fieldM) are driven by
mid\_mem, which reflects L2's hidden state under \emph{random} waves.
Random waves inject no zone structure into mid\_mem; each update partially
overwrites the birth-seeded structure.
Over $\sim 6$ collapses per site (3000 steps), the injected structure cannot
survive ongoing erasure by the unstructured update signal.

\textit{What ctrl does right.}
With random waves, ctrl's copy-forward loop accumulates internally consistent
spatial structure via drift from its own (initially noisy) fieldM.
Over 3000 steps this drift produces modest but genuine zone-like encoding
(na\_ratio$=1.245 > 1$) without any external coordination.

The efficient relay primitive must impose zone-structured viability events AND
a matching zone-convention birth seed in the receiving module simultaneously.
This is a harder design problem than either geo or mean\_relay separately,
and constitutes the open question answered definitively in Paper~50.

\subsection{Zone identity encoding is not transmissible (Paper 50)}

Paper~50 audits the relay result from Paper~49 across five conditions
to identify the mechanism behind geo's parity encoding and test whether
any relay coupling recovers identity encoding in L2.

\textbf{The only condition that achieves identity encoding is ctrl\_std.}
When L2 runs its own WaveEnvStd (same zone parameters as L1, independent
random seed), it achieves na\_ratio$=1.219$, matching L1's $1.247$.
This confirms that identity encoding is not special to L1: it is a
property of the VCML formation process under structured zone-launched waves.

\textbf{No relay coupling achieves identity encoding.}
\begin{itemize}
  \item geo\_copy (instantaneous amplitude copy): na$=0.778$ [parity]
  \item geo\_4class (4-level amplitude grading): na$=0.949$ [parity, improved]
  \item geo\_own\_relay (own WaveEnvStd + L1 birth seeding): na$=0.832$ [parity]
\end{itemize}

\textbf{Birth seeding from L1 actively degrades ctrl\_std.}
geo\_own\_relay --- L2's own WaveEnvStd plus L1 zone-mean birth seeding ---
achieves na$=0.832$, substantially worse than ctrl\_std's $1.219$.
L1 and L2 develop different but equally valid identity-encoding patterns
(different spatial orientations in HS$=2$ space).
L1's birth seeds conflict with L2's self-consistently forming pattern:
\emph{pattern-mismatch interference}, not a signal quality problem.

\textit{Why identity cannot be relayed.}
L1's identity encoding accumulates over 3000 steps of wave events from
zone-anchored launch positions, reinforced at each step by the copy-forward
loop ($\mathrm{hid}_\mathrm{new} \approx \beta_s \cdot \mathrm{fieldM}[\mathrm{cell}]$).
An instantaneous snapshot of L1's wave amplitude array loses:
(1)~temporal integration history --- the spatial gradient lives in L1's fieldM,
not in any single wave amplitude snapshot;
(2)~copy-forward self-consistency --- L2's fieldM under geo relay is pulled toward
parity by the viability signal and toward identity by its own copy-forward loop,
but the parity viability signal dominates.

\textit{Minimum requirement for identity encoding.}
L2 must run an independent WaveEnvStd with zone-anchored launch positions.
No relay shortcut achieves this.
This rules out all tested relay architectures as mechanisms for identity-preserving
inter-module communication.

\textit{Caveat on Papers 35--39.}
The relay gains of $3$--$5\times$ in sg4 reported in Papers~35--39 are real and
correctly measured.
However, in light of Paper~50, the structure is parity encoding (na$<1$),
not identity encoding (na$>1$).
``Zone structure'' in relay contexts should be read as zone amplitude separation
(sg4), not zone topographic identity (na\_ratio).
Chain stability, cleft robustness, and depth-stability results stand;
the interpretation of what structure is being relayed requires this caveat.

\subsection{The research process as VCSM instance (Paper 51)}

Paper~51 observes that the methodology that produced Papers~1--50 is
structurally identical to the VCSM framework.
The research process operates on three timescales:
session (context window, analogous to baseline\_h),
paper (MEMORY.md accumulation, analogous to mid\_mem),
and framework document (VCML\_THEORY\_COMPLETE.tex, analogous to fieldM).
Theoretical consolidation is gated on adversarial stress testing
(the quiescence condition: findings must survive ablation and metric challenge
before entering the framework document).
Every paper terminates when its question closes (mandatory turnover).
Each new paper reads the accumulated theory and memory files (copy-forward loop).

The consequence is that the framework achieves zone identity encoding ---
topographic, gradient-structured theory --- by running the full formation process,
not by receiving it from a prior source.
Zone identity encoding in the theory is not transmissible (Paper~50):
reading the framework document produces parity understanding (what is
suppressed vs.\ excited) but not identity understanding (the topographic gradient
of why each primitive survived, which constraints forced which mechanisms).
That gradient lives in the formation history: the 50-paper viability-gated ladder,
the 11 null results and their specific failure modes, the metric evolution
from sg4 to na\_ratio, the adversarial challenges and what they broke.

\textit{Centerpiece example.}
The parity/identity distinction was invisible in sg4 for approximately
Papers~35--49 ($\sim$15 papers).
It became visible in one step when na\_ratio was derived from the theoretical
primitive (zone topography requires non-adjacent zones to be more distinct
than adjacent zones) rather than from measurement convenience.
The relay architecture appeared successful for 15 papers because the question
being asked (are zones separated?) was not the question the theory requires
(are the \emph{correct} zones separated?).
The correct question was invisible until the theory, built by constraint
propagation, outgrew the metric.

\subsection{Two-channel write theory (Paper 52)}
\label{subsec:paper52}

Paper~52 asks whether the \emph{signed direction} of the VCSM contrast signal
($\mathrm{mid\_mem} = \mathrm{hid} - \mathrm{baseline\_h}$) is load-bearing for
zone identity encoding, or whether its magnitude alone suffices.
Nine consolidation-write-policy conditions were tested in a two-phase protocol:
2000~steps of formation (standard WaveEnvStd) followed by 1000~steps of adversarial
environment (WaveEnvFlipped: suppress/excite convention inverted).

\textit{C\_abs is the only parity failure.}
Taking the absolute value of mid\_mem (condition C\_abs) is the unique condition that
fails identity encoding (na\_ratio = 0.877, parity territory).
All eight signed conditions achieve identity encoding (na\_ratio = 1.03--1.80).
The signed direction of the contrast signal encodes which side of the viability
boundary the cell occupies: moving toward collapse (negative) versus toward survival
(positive).
Collapsing this to $|\mathrm{mid}|$ destroys the directional information that makes
intergenerational propagation zone-informative.
The boundary does not generate information; the \emph{sign of the crossing} does.

\textit{Two-channel write theory.}
Paper~52 identifies two independent axes governing the write-policy space:

\begin{enumerate}
  \item \textbf{Content axis} (what to write):
    C\_raw (raw hid, na=1.802) $>$ signed mid ($\sim$1.2--1.3) $>$ $|$mid$|$ (0.877).
    The raw hidden state carries stronger zone identity signal than mid\_mem;
    subtracting baseline\_h removes some identity information while adding
    boundary sensitivity.
  \item \textbf{Selectivity axis} (when to write):
    C\_perturb (cos\_maint=0.664, writes-per-site=0.0004) $>$ C\_near (0.546) $\gg$
    C\_ref ($-0.274$).
    Write \emph{trigger policy} --- not write content --- determines adversarial resistance.
    C\_perturb writes only during active perturbation events, achieving best adversarial
    maintenance at 83$\times$ fewer writes than C\_near.
\end{enumerate}

\textit{C\_ref and C\_far re-encode.}
Both conditions write at high frequency and produce negative cosine maintenance
(re-encoding the adversarial environment rather than preserving formation structure).
High write frequency = rapid adaptation = inability to distinguish formation from
adversarial phase.
All near-boundary conditions preserve identity under adversarial (na\_ratio $> 1$
after flip), confirming that boundary-gated selectivity confers structural stability.

\textit{Theory implication.}
Standard VCSM (C\_ref) lies in the responsive quadrant of the two-channel write space:
high selectivity on content axis (signed mid over raw hid), low selectivity on
trigger axis (all quiescent cells write).
The two-channel space defines a design tradeoff:
moving along the content axis (toward C\_raw) increases formation encoding strength;
moving along the selectivity axis (toward C\_perturb or C\_near) increases adversarial
resistance at the cost of slower formation.
C\_rand (random trigger, matched write rate) achieves na=1.335 --- the random trigger
selects a different valid subsample of the formation signal, confirming that
boundary proximity is one valid selection criterion, not the only one.

\subsection{The VCML behavioral phase diagram (Paper 53)}
\label{subsec:paper53}

Paper~53 synthesizes the experimental evidence from Papers~0--52 into a
unified behavioral phase diagram for VCML.
The central finding is that each experimental surprise corresponded to a regime
crossing along an axis that was not yet separated in the mental model.
Five independent control axes are identified.

\textit{Axis~I: Perturbation intensity (coll/site).}
sg4 is non-monotone: peak at coll/site~$\approx 0.004$,
falling in both the sub-optimal (formation-starved) and over-perturbed
(consolidation-disrupted) regimes.
The peak is the point where both the copy-forward event rate and the
quiescence streak accumulation are simultaneously satisfied.

\textit{Axis~II: Spatial resolution ($\delta = W_\text{zone}/r_\text{wave}$).}
Hard threshold at $\delta \approx 2.5$.
Below this, waves interpenetrate adjacent zones (wave bleed), producing parity
encoding or no structure.
The transition is categorical (hard threshold from wave-bleed physics), not
smooth.

\textit{Axis~III: Formation flux ($\Omega = \mathrm{WR} \times \phi_w$).}
The operative coordinate for perturbation rate (Papers~46--47).
Coupled to Axis~I through WR: the optimal WR is a function of grid size, not
a single universal constant.

\textit{Axis~IV: Write policy (2D).}
The content axis (C\_raw $>$ signed mid $>$ $|$mid$|$, identity encoding
strength) and the trigger selectivity axis (C\_perturb $>$ C\_near $>$ C\_ref,
adversarial resistance) are orthogonal.
These axes interact only at the C\_abs extreme, where unsigned content and
boundary trigger combine to produce the unique parity failure.

\textit{Axis~V: Scale and finite correlation length.}
The copy-forward loop has correlation length $\sim r_\text{wave} \approx 2$
sites.
Above $\sim$25,000 sites per zone, zone-wide gradients cannot be maintained
and structure forms in patches.
The dynamic advantage is preserved at all scales
($+81\%$ to $+829\%$, Papers~V92--V94), but the spatial organization
degrades.

\textit{Cross-axis coupling.}
Axes~I and~III are coupled through WR (increasing WR changes both).
Axes~II and~V interact through zone width (large N increases $\delta$ without
fixing patch formation, which is a distinct correlation-length mechanism).
Axis~IV (C\_perturb) is contingent on Axis~I being in the adaptive zone:
C\_perturb writes nothing when perturbation events are absent.

\textit{VCML as a minimal causal field system.}
A minimal causal field system generates, from the same local rules in different
parameter regimes: identity encoding, adversarial resistance, intergenerational
memory, and spatial topology.
VCML satisfies all four.
The phase diagram is the evidence: the same system in the over-perturbed regime
loses identity encoding while retaining turnover; in the C\_abs regime, it
retains adversarial stability while losing identity encoding; at S4, it retains
formation while losing zone-wide coherence.
The correct question is not whether the mechanism is understood, but which
regime of the phase diagram a given experiment inhabits.

\textit{Paper~53 revision (Paper~54):}
The cross-axis prediction that C\_perturb fails at low amplitude was inverted
by experiment.
C\_perturb maintains identity encoding at all amplitude levels tested
(na\_ratio = 1.49, 1.39, 1.25 for sub-threshold, optimal, and over-perturbed),
while C\_ref fails at both extremes.
This corrects the mechanism: Axis~IV trigger selectivity is \emph{noise-gating}
(minimizes $\sigma_w$), not signal-boosting.
The $\sim$10--20$\times$ lower within-zone variance of C\_perturb produces
high SNR even with small absolute sg4.

\subsection{Phase diagram refinements (Paper 54)}
\label{subsec:paper54}

Paper~54 tests three open questions from Paper~53 and revises two items.

\textit{Order parameter.}
sg\_C = sg4/$\sigma_w$ is not the unified order parameter.
snr is flat ($\approx 0.105$) across the identity regime and drops slightly in
the parity regime, providing no additional discriminating power over na\_ratio.
The candidate order parameter is a vectorial SNR:
$C_\text{order} = \|\mu_\text{nonadj} - \mu_\text{adj}\| / \sigma_w$,
measuring whether correct zone pairs are separated beyond within-zone noise.

\textit{Correlation length law.}
At S2 (zone\_width=20), sweeping $r_\text{wave}$ from 1 to 16 reveals two
distinct failure modes.
At $r_\text{wave}=1$ ($\delta=20$, above the wave-bleed threshold),
na\_ratio = 0.748 (parity): a correlation-length failure independent of
$\delta$.
Empirical law: identity encoding requires
zone\_width $\leq 5 r_\text{wave}$.
This is separable from $\delta$: both conditions (zone containment by $\delta$,
loop coherence by $r_\text{wave}$) must be satisfied simultaneously.

\textit{Trigger mechanism revision.}
The Paper~53 claim that Axis~IV trigger selectivity selects more informative
writes is incorrect.
Trigger selectivity is a noise-gating operation: C\_perturb writes to fewer
cells, keeping $\sigma_w$ near zero, producing high SNR from tiny sg4.
This mechanism accounts for both C\_perturb's adversarial resistance
(Paper~52) and its cross-amplitude robustness (Paper~54).

\subsection{Lyapunov functional for identity formation (Paper 55)}

Paper~55 tests whether C\_order has monotone positive drift (Lyapunov property)
by tracking it over full 3000-step trajectories across five behavioral regimes.

\textbf{Result:}
C\_order shows monotone positive drift \emph{only} in the noise-gated condition
(C\_perturb): C\_order rises from $+0.011$ at step~1000 to $+0.262$ at step~2000
and plateaus at $+0.159$ at step~3000.
All C\_ref conditions produce C\_order $\approx 0$ or negative throughout,
because $\sigma_w \approx 0.15$--$0.32$ overwhelms the small inter-zone contrast.

\textbf{Transient vs.\ stable identity.}
The standard C\_ref optimal condition (supp=0.25, exc=0.50) reaches na\_ratio = 1.026
at step~2000 (the measurement point of Paper~54), but C\_order $= -0.013$ at
this same timepoint.
By step~3000, na\_ratio falls to 0.729 (parity).
The Paper~54 measurement captured a transient identity peak, not a stable attractor.
C\_perturb identity is stable: na = 1.412 and C\_order = $+0.159$ at step~3000.

\textbf{Lyapunov condition.}
C\_order $= (D_\mathrm{nonadj} - D_\mathrm{adj})/\sigma_w$.
Let $S = D_\mathrm{nonadj} - D_\mathrm{adj}$ and $W = \sigma_w$:
\begin{equation}
\dot{C}_\mathrm{order} > 0 \quad\Longleftrightarrow\quad \dot{S}/S > \dot{W}/W
\end{equation}
C\_perturb satisfies this because $\dot{W}/W \approx 0$ (trigger selectivity
limits writes, stabilising $\sigma_w$).
C\_ref violates it because $\dot{W}/W > 0$ (writes to all quiescent cells
continuously accumulate noise).

\textbf{Three jointly necessary conditions} for the Lyapunov property:
(1)~zone\_width $\leq 5r_\text{wave}$ (correlation length, ensures copy-forward
self-amplification $K_\mathrm{self} > 0$);
(2)~perturbation amplitude in the identity window (Axis~I);
(3)~$\dot{W}/W \approx 0$ (noise floor control).
All three correspond to regime conditions already identified in Papers~53--54.
The Lyapunov analysis does not discover new boundaries; it explains why they matter
at the level of dynamical stability.

\subsection{Three completions for the regime diagram (Paper 56)}

Paper~56 closes three questions left open by Paper~55:

\textbf{(A) SS threshold cannot substitute for C\_perturb noise-gating.}
Sweeping the consolidation threshold SS over $\{10, 20, 30, 50, 100, 200\}$
in the sub-threshold C\_ref condition shows that SS$=30$ produces marginal
identity (na$=1.014$), but no level achieves the Lyapunov property
(cross\_step $=$ never for all SS).
Signal attenuation by the mid-memory decay ($\mathrm{MID\_DECAY}^{SS}
= 0.99^{200} = 0.134$) cancels the write-frequency reduction.
Higher SS lowers write\_ready\_frac (0.732 at SS$=10$ to 0.536 at SS$=200$)
but proportionally weakens the consolidated signal; the Lyapunov condition
remains unsatisfied.

\textbf{(B) C\_perturb extends the effective correlation-length threshold.}
Paper~54 established identity requires zone\_width $\leq 5r_\text{wave}$
(critical $\delta \leq 5$) for C\_ref.
Under C\_perturb at the S2 grid, $r_\text{wave}=1$ achieves identity
(na~$=1.160$, $C_\text{order}=+0.041$) at $\delta=20$, and $r_\text{wave}=2$
reaches na~$=1.261$ at $\delta=10$.
The effective threshold extends at least $2$--$4\times$ beyond the C\_ref
value.
Mechanism: noise-gated writes keep $\sigma_w \approx 0.001$--$0.011$
vs.\ C\_ref $\sigma_w \approx 0.15$--$0.19$; the elevated SNR compensates for
small wave radius.
The optimal $r_\text{wave}$ under C\_perturb is $r=2$; $r=8$ fails because
large footprints cause wave-class overlap, inflating $\sigma_w$ toward the C\_ref
range.

\textbf{(C) na\_ratio is write-frequency-invariant.}
write\_ready\_frac (instantaneous fraction of cells passing the consolidation
gate) ranges from $0.31$ to $0.96$ across all conditions with no monotone
relationship to na\_ratio.
C\_perturb achieves identity at $r \in \{1,2,4\}$ despite operating in a very
different write-frequency regime than any C\_ref condition.
Content selectivity (which events trigger writes) dominates write frequency
in determining identity formation.
Corollary: na\_ratio~$-1$ is the robust, write-policy-agnostic order parameter
for VCML identity; $C_\text{order}$ provides mechanistic resolution but is not
comparable across write policies.

\subsection{Analytic closure of the Lyapunov condition (Paper 57)}

Paper~57 derives the noise-growth term $\dot{W}/W$ analytically from wave geometry,
closing the derivation left open in Paper~55.

\textbf{Geometric approximation.}
Cells within $r_\text{wave}$ of a zone boundary receive mixed-zone wave input.
The boundary-contaminated fraction is $f_\text{bnd} = 2r_\text{wave}/W_\text{zone} = 2/\delta$;
the interior fraction is $f_\text{int}(\delta) = (1-2/\delta)^+$.
Under C\_perturb, interior-cell writes dominate; under C\_ref, boundary cells also
write, giving $\mathbb{E}[\Delta\sigma_w^2 | \text{C\_ref}] \approx (2/\delta)\,V_\text{bnd} > 0$
at all finite $\delta$.
Critical threshold: $\delta_c = 2$.

\textbf{Experimental verification (Exp~57).}
Fine $r_\text{wave}$ sweep $\{2,3,4,5,6,7,8,10\}$ under C\_perturb, S2 sub-threshold,
40 runs: $\sigma_w$ rises monotonically with $r_\text{wave}$, confirming the
boundary-contamination model.
Deviation at $r_\text{wave}=10$ ($\delta=2$, $f_\text{int}=0$): geometry predicts
full collapse but system achieves na~$=1.091$.
Explanation: same-zone launch area asymmetry maintains positive causal purity
$P_\text{causal}$ even when $f_\text{int}=0$.

\textbf{Substrate-agnostic form.}
The correct primitive is \emph{causal purity} $P_\text{causal}$: the fraction of
pre-consolidation wave-contact events from the same zone class.
The Lyapunov condition becomes $\dot{C}_\text{order} > 0 \iff P_\text{causal} > p_c$
for threshold $p_c \approx 1/3$.
This statement contains no lattice-specific quantities.
It unifies the three prior laws (Papers~54--56): correlation-length threshold, noise-floor
condition, and new geometric threshold are all projections of $P_\text{causal} > p_c$.

\subsection{Two-dimensional causal purity (Paper 58)}

Paper~58 implements direct empirical measurement of $P_\text{causal}$: for every
consolidation event, the zone index of the triggering wave is recorded and compared
to the cell's zone. No geometric approximation.

\textbf{Finding 1.} Empirical spatial purity $P_\text{sp}$ (0.88--0.98) matches the
geometric prediction and decreases with $r_\text{wave}$.

\textbf{Finding 2.} C\_perturb and C\_ref have nearly identical $P_\text{sp}$
(within 0.02--0.03 at each $r_\text{wave}$), yet $\sigma_w$ differs 3--10$\times$.
Spatial causal purity alone cannot explain the write-policy gap.

\textbf{Finding 3.} The linear noise model $\sigma_w \approx A(1-P_\text{sp})$
(with $A \approx 0.84$) holds for C\_perturb across all $r_\text{wave}$, but fails
for C\_ref (ratio 1.6--10$\times$ above linear).
The missing dimension is \emph{temporal causal purity} $P_\text{temp}$: the fraction
of consolidation events that occur \emph{during} active causal-wave contact.
C\_perturb has $P_\text{temp}=1$ by construction; C\_ref has $P_\text{temp} < 1$
(writes during quiescence, reading stale, decayed mid-memory).

\textbf{Two-dimensional identity condition.}
\begin{equation*}
  \dot{C}_\text{order} > 0 \iff P_\text{sp} > p_c^\text{sp}
  \;\text{ AND }\; P_\text{temp} > p_c^\text{temp}
\end{equation*}
Extended noise model: $\sigma_w \approx A_\text{CP}(1-P_\text{sp}) + B(1-P_\text{temp})$
with $B \gg A_\text{CP}$.

\textbf{Primitive restated.}
The substrate-agnostic primitive is \emph{distinction + causal propagation + persistence}.
Carrier mortality is the mechanism enforcing the persistence requirement on biological
substrates, not a primitive requirement in itself.
Both purity conditions ($P_\text{sp}$, $P_\text{temp}$) are requirements on the
persistence of a causally attributed distinction under substrate change.

\subsection{Field equation of VCML and the r\_wave=8 anomaly (Paper 59)}

Paper~59 hypothesises a third purity dimension, \emph{historical purity}
$P_\text{hist}$ (fraction of mid\_mem content at consolidation from same-zone waves),
and measures it empirically.

\textbf{Null result.}
$P_\text{hist} \approx P_\text{sp}$ within $0.04$ across all conditions.
The geometric prediction for $P_\text{hist}$ is severely wrong
(predicts 0.68--0.84; actual 0.93--0.98), corrected by the same launch-area asymmetry
from Paper~57.
$P_\text{hist}$ is not a distinct third dimension.

\textbf{Model A field equation.}
The continuous-time limit of the GRU--CML--fieldM update is:
\begin{equation*}
  \partial_t \phi(x,t)
  = -\kappa\,\phi + D\nabla^2\phi + r_w\,\mathrm{FA}\,s(x) + \eta(x,t),
\end{equation*}
a Model~A (Ornstein--Uhlenbeck) field theory with noise amplitude
$\langle\eta^2\rangle \propto (1-P_\text{sp})^2 V_\text{bnd} + (1-P_\text{temp})^2 V_\text{temp}$.
$C_\text{order} = S/\sigma_w$ is the field-theoretic SNR
(order-parameter amplitude / fluctuation amplitude).

\textbf{$r_\mathrm{wave}=8$ parity-resonance failure.}
The anomaly is a structural (not purity) failure.
At $r_\mathrm{wave} \approx W_\mathrm{zone}/2.5$, adjacent zones exchange
excitatory/inhibitory polarity signal maximally, but waves do not yet reach
non-adjacent zones.
fieldM encodes zone polarity (even = suppressed, odd = excited) rather than
zone position, inverting the na signature ($D_\text{adj} > D_\text{nonadj}$).
Third condition: $r_\mathrm{wave} < W_\mathrm{zone}/3$ OR $r_\mathrm{wave} > W_\mathrm{zone}/2$
(drive-pattern non-resonance), orthogonal to causal purity.

% ─────────────────────────────────────────────────────────────────────────────
\section{The Parameter Manifold: Why No Single Equation Suffices}
\label{sec:manifold}

\subsection{The 8D product structure}

The VCML system lives in an 8-dimensional product manifold $C_1 \times C_2$,
where the two factors govern qualitatively different aspects of behaviour.

\medskip
\textbf{$C_1$ --- Viability volume} $(\nu,\, \mathrm{FD},\, \mathrm{FA},\, \beta_s)$.
These four parameters determine \emph{whether} spatial memory consolidation occurs at all.
$\nu$ sets the turnover rate; FD (field decay) controls how fast unconsolidated
structure is lost; FA (force amplitude) sets the signal entering mid-memory;
$\beta_s$ (birth-seed fraction) determines how much of the parent fieldM is
inherited at birth.
The adaptive regime occupies the interior of $C_1$; the boundaries lead to either
the crystallised regime (too-slow turnover, fieldM locks) or the sub-viability
regime (too-fast, nothing consolidates).

\medskip
\textbf{$C_2$ --- Structural manifold} $(\kappa,\, \nu/\kappa,\, \beta_s,\, W_\mathrm{zone})$.
These four parameters determine \emph{what kind} of spatial structure forms,
given that consolidation is occurring.
$\kappa$ (diffusion constant) sets the smoothing length; $\nu/\kappa$ is the
dimensionless ratio that selects the universality class (Allen-Cahn at $\nu/\kappa \ll 1$,
Burgers/KPZ at $\nu/\kappa \sim 1$); $W_\mathrm{zone}$ is the zone width relative
to the wave correlation length.

\medskip
\textbf{Causal reparametrisation of $C_2$.}
The lattice coordinates $(\kappa,\, \nu/\kappa,\, W_\mathrm{zone})$ are valid for the
lattice substrate but reference Euclidean geometry explicitly.
Paper~36 established that the operative correlation length is $\xi \approx r_w$
(causal propagation radius), not $\sqrt{\kappa/\nu}$ (diffusion length).
This motivates a substrate-agnostic reparametrisation:

\begin{center}
\begin{tabular}{lll}
\toprule
\textbf{Lattice parameter} & \textbf{Causal coordinate} & \textbf{Definition} \\
\midrule
$\kappa$ & wave fanout $\phi_w$ & cells reached by one wave event $= \pi r_w^2$ on lattice \\
$\nu/\kappa$ & $\Omega = \nu \cdot \phi_w \cdot T_{1/2}$ & collapse events per causal neighbourhood per $T_{1/2}$ \\
$W_\mathrm{zone}$ & causal depth ratio $\delta = W_\mathrm{zone}/r_w$ & zone width in units of causal horizon \\
$\beta_s$ & $\beta_s$ (unchanged) & copy-forward inheritance fraction \\
\bottomrule
\end{tabular}
\end{center}

$\Omega \gg 1$: rapid copy-forward refresh $\to$ Burgers/sharp-boundary regime.
$\Omega \ll 1$: slow decay $\to$ Allen-Cahn/diffuse regime.
$\delta \lesssim 1$ (S1): zone fits inside one causal neighbourhood; zone-wide coherence possible.
$\delta \gg 1$ (S4): independent causal neighbourhoods; structure forms in patches of size $\sim r_w$ (V94).
The relay threshold $N_\mathrm{crit} = L/(5r_w)$ is $\delta \geq 5$: five causal neighbourhoods
required for a reliable zone-mean estimate (Papers~35--37).

On a graph substrate, $\phi_w$ becomes the mean causal neighbourhood size under the graph topology
(not $\pi r_w^2$), $\delta$ becomes zone nodes / neighbourhood size, and $\Omega$ uses
graph-neighbourhood $\phi_w$ directly.
The reparametrisation eliminates the Euclidean assumption from $C_2$ without changing
any experimental result: every lattice result is reproduced by substituting
$\phi_w = \pi r_w^2$, $\delta = W_\mathrm{zone}/r_w$.
One additional parameter appears on general graphs:
birth fanout $\phi_b$ (cells reachable by copy-forward seeding in one step);
on the lattice $\phi_b = 1$ (nearest-neighbour, degenerate) and cannot be measured independently.

\medskip
\textbf{The shared axis.}
$\beta_s$ is the \emph{only} parameter appearing in both manifolds.
It couples viability (how much structure survives birth) to spatial structure
(how coherently that structure propagates).
Because the coupling runs through a single scalar, $C_1$ and $C_2$ can be
explored and optimised largely independently --- a fact exploited throughout
Papers~7--34 to close each axis in isolation.

\subsection{The closest unified description, and why it is still a slice}

The stochastic Allen-Cahn/KPZ PDE
\begin{equation}
\frac{\partial F}{\partial t}
  = \underbrace{\kappa\,\nabla^2 F}_{\text{diffusion}}
  - \underbrace{\lambda(F - F^*)^2}_{\text{Allen-Cahn attraction}}
  + \underbrace{\tfrac{1}{2}\chi(\nabla F)^2}_{\text{KPZ/Burgers}}
  + \underbrace{\eta(x,t)}_{\text{stochastic drive}}
\label{eq:unified}
\end{equation}
is the closest single equation to a unified description of within-zone fieldM dynamics.
With $\chi \to 0$ it reduces to Allen-Cahn (smooth-interface, $\nu/\kappa \ll 1$ regime);
with $\lambda \to 0$ it reduces to KPZ (rough-interface, $\nu/\kappa \sim 1$ regime).
The crossover is controlled by $\nu/\kappa$ (Paper~24).

Yet Eq.~(\ref{eq:unified}) is \emph{still only a slice}.
It lives entirely in $C_2$: it says nothing about whether the adaptive window is
open ($C_1$), how fast zone means grow (ODE closure), when the field commits
($t^*$ scaling, Paper~27), or whether two coupled regions can share structure
(relay conditions, Papers~35--37).
The full theory requires a description of the product $C_1 \times C_2 \cong \mathbb{R}^8$,
and no single PDE captures all eight dimensions.

\subsection{Each equation is a projection}

The system genuinely lives on the product surface.
Each equation in the theory is the \emph{projection} of that 8D surface onto the
slice relevant to a particular question:

\begin{center}
\begin{tabular}{lll}
\toprule
\textbf{Question} & \textbf{Relevant slice} & \textbf{Dominant equation} \\
\midrule
Does structure form at all?       & $C_1$: $(\nu$, FA, FD$)$              & Adaptive window bounds \\
What causal structure emerges?    & $C_2$: $(\phi_w$, $\Omega$, $\delta)$ & Allen-Cahn ($\Omega\!\ll\!1$) / KPZ ($\Omega\!\gg\!1$) \\
How fast does fieldM grow?        & Zone-mean projection ($\beta_s$, $P_c$) & ODE closure \\
When does the system commit?      & $P_c \cdot \mathrm{FA} \cdot \mathrm{WR}$ & $t^*$ scaling \\
Can two regions share structure?  & Relay geometry ($r_w$, $L$)           & $N < L/(5r_w)$, $\Phi < 1$ \\
Is memory maintained after turnover? & $C_1 \cap C_2$: gate + $\beta_s$   & Copy-forward loop \\
\bottomrule
\end{tabular}
\end{center}

In any given parameter regime one term dominates and one equation suffices ---
exactly as QFT and general relativity each describe their domain correctly
with the full theory requiring both.
This is not a weakness of the theory; it is a consequence of the geometry.
A single equation would be a false unification: it would collapse the $C_1$--$C_2$
distinction and hide the fact that viability and structure are controlled by
different physical mechanisms with different experimental handles.

% ─────────────────────────────────────────────────────────────────────────────
\section{The Complete Parameter Space}

\subsection{Key equations}

\begin{align}
\text{(ODE closure)} \quad &\frac{d\bar{F}}{dt} = P_c \cdot \mathrm{FA} \cdot (\bar{m}-\bar{F}) - \nu\bar{F} \\[4pt]
\text{(Saturation)} \quad  &\mathrm{sg4}(\mathrm{FA}) = \frac{C \cdot \mathrm{FA}}{\mathrm{FA} + K_\mathrm{eff}}, \quad K_\mathrm{eff} \approx 0.114 \\[4pt]
\text{(Spatial amplif.)} \quad &\Gamma \sim W_\mathrm{zone}^{0.6}\,\kappa^{-0.3}, \quad K_\mathrm{eff}^\mathrm{spatial} \approx 0.040 \\[4pt]
\text{(Zone width)} \quad  &W_\mathrm{zone}^* \sim \sqrt{\kappa/\nu} \\[4pt]
\text{(Commit epoch)} \quad &t^* \sim 1/(P_c \cdot \mathrm{FA} \cdot \mathrm{WR}) \\[4pt]
\text{(Relay)} \quad        &N < L/(5\,r_w) \;\text{ and }\; \Phi = \mathrm{WR}\,r_w/T_\mathrm{dur} < 1 \\[4pt]
\text{(Non-gradient)} \quad &\mathrm{curl}(f_F, f_m) = -\mathrm{FA} \neq 0
\end{align}

\subsection{Hierarchy of descriptions}

\begin{center}
\begin{tabular}{lll}
\toprule
Scale & Description & Papers \\
\midrule
Cell          & 5-step VCSM rule                    & 0 \\
Within-zone   & Allen-Cahn/Burgers stochastic PDE   & 22--26 \\
Zone-mean     & ODE: $d\bar{F}/dt = P_c\,\mathrm{FA}(\bar{m}-\bar{F}) - \nu\bar{F}$ & 20, 23 \\
Parameter     & 4D viability volume $C_1$            & 7--13 \\
Structure     & 4D structural manifold $C_2$         & 22--26, 30 \\
Multi-region  & Geographic relay, cleft geometry    & 35--39 \\
Dynamics      & Metastability, CL, copy-forward refresh & 38--42 \\
Robustness    & Rule-family robustness, formation vs.\ maintenance & 43 \\
Adversarial   & Imprinting critique closed, memory vs.\ environment  & 44 \\
Formal        & Non-gradient (Lean 4 verified)      & 33 \\
\bottomrule
\end{tabular}
\end{center}

% ─────────────────────────────────────────────────────────────────────────────
\section{Biological Predictions}

All predictions follow directly from the theory; none require additional assumptions.

\paragraph{Adult hippocampal neurogenesis.}
VCML death--birth cycle $\equiv$ adult neurogenesis in the dentate gyrus.
The copy-forward mechanism proposes that newborn granule cells inherit structured fieldM
from mature neighbours; dendritic contact during migration is the candidate substrate for
this inheritance step (\textit{mechanistic speculation --- not established}).

\textit{What is confirmed in animal models:} preventing apoptosis of adult-generated DG
neurons impairs pattern separation.
This is consistent with VCML: crystallised dynamics (no death) degrades zone structure
maintenance relative to the adaptive regime.

\textit{VCML-specific prediction (not yet tested):} blocking adult neurogenesis should
impair \emph{maintained} spatial memory (zone structure persistence after input removal)
more than \emph{initial acquisition} (zone formation).
This distinguishes the maintenance mechanism from a simple capacity account and is a
different experiment from the standard neurogenesis--learning literature.
The mechanism is not ``more neurons = more capacity'' but ``optimal turnover rate maintains
adaptive zone structure via copy-forward refreshing''.

\paragraph{Sleep replay.}
Offline replay of multiple experiences during sleep $\equiv$ Paper~41's task switching.
\textbf{Prediction:} replay maintains hippocampal zone structure $5$--$10\times$ better than
single-trace consolidation alone.
The optimum replay rate $\approx 0.6 \times \tau_{1/2}$ (where $\tau_{1/2}$ is the
zone-structure half-life without replay).
Absence of replay causes structural decay --- not just poor encoding, but loss of the spatial
index itself.

\paragraph{Topographic map formation.}
Geographic relay coupling $\equiv$ axonal topographic projection between cortical areas.
\textbf{Prediction:} deterministic topographic misalignment (consistent offset) does not
impair downstream zone formation; stochastic scrambling of terminal fields does.
Critical parameter: $\Phi = \mathrm{WR}\,r_w/T_\mathrm{dur} < 1$ for relay to succeed.
This maps to a condition on the ratio of axonal conduction velocity to synaptic time
constant in the biological regime.

\paragraph{Hippocampal vs.\ cortical memory.}
High-turnover regions (hippocampus, $\nu$ large) operate in the adaptive regime and
exhibit gradual forgetting on the metastability timescale.
Low-turnover regions (cortex, $\nu \approx 0$) operate near the crystallised regime and
exhibit near-permanent memory at the cost of reduced plasticity.
The dentate gyrus sits at the adaptive--crystallised boundary by design: high enough $\nu$
to maintain zone structure via copy-forward; low enough to not decay before consolidation.

% ─────────────────────────────────────────────────────────────────────────────
\section{State of the Theory}

\subsection{Single-region theory: complete mechanistic account}

As of Paper~44, the dominant mechanism underlying single-region VCML dynamics is
established, rule-family robust, and adversarially validated.

Paper~42 resolved all three open anomalies (C$_\mathrm{loc}$ ratio$>1$, C$_\mathrm{half}$
surge, T$_\mathrm{sw}$ reversal).
Paper~43 demonstrated robustness: formation persists across all rule variants and ablations;
maintenance requires viability-gated copy-forward.
Paper~27 provided predictive validation: the commitment timescale scaling
$t^* \sim 1/(P_c \cdot \mathrm{FA} \cdot \mathrm{WR})$ was predicted then confirmed.
Paper~44 closed the imprinting critique: the adaptive mechanism maintains positive cosine
fidelity with the original encoding under 2000 steps of sustained adversarial input,
while the passive (no-gating) mechanism encodes the adversarial pattern
($\Delta = 0.205$, $p < 0.001$).

The dominant single-region phenomenology is captured by one sentence:

\begin{quote}
\textit{The copy-forward loop (viability-gated consolidation $\to$ structured collapse $\to$
birth seeded from fieldM) maintains zone structure, and its rate is determined by the
structured-collapse rate from any perturbation source (waves, task switching, boundary
mismatch, or intrinsic noise); structure formation is additionally robust to rule variants
because zone-differentiated wave input is sufficient to drive differentiation under any
consolidation mechanism.}
\end{quote}

\textbf{What this claim does and does not assert:}
\begin{itemize}[itemsep=2pt]
\item \textit{Does assert:} viability-gated copy-forward is the dominant maintenance mechanism;
  it is necessary (removing gating degrades maintenance to near field-decay-only);
  it is robust (six $\pm50\%$ rule variants all maintain structure).
\item \textit{Does not assert:} this is the \emph{unique} possible mechanism for formation
  (any local rule with zone-differentiated input likely produces formation);
  the ODE/PDE are exact (they are leading-order approximations valid near the mean field);
  the theory is mathematically complete in the sense of formal proof theory.
\end{itemize}

\subsection{Open frontiers}

\textbf{Finite causal horizon and trans-horizon coherence.}
Paper~36 established that the effective correlation length is $\xi \approx r_w$ ---
determined by the causal propagation radius of wave events, not by diffusion length $\sqrt{D}$
or any $C_1$/$C_2$ parameter.
This implies a \emph{finite causal horizon}: two cells separated by $d > r_w$ cannot share
a common wave ancestor in a single step and therefore cannot be directly causally correlated.

Two regimes follow directly from this bound:
\begin{itemize}[itemsep=1pt]
  \item \emph{Sub-horizon zones} ($W_\mathrm{zone} \lesssim r_w$): the entire zone lies within
    one causal neighbourhood; zone-wide coherence is achievable.
  \item \emph{Trans-horizon zones} ($W_\mathrm{zone} \gg r_w$): a single zone contains
    $(W_\mathrm{zone}/r_w)^2$ independent causal neighbourhoods; structure forms in patches
    of size $\sim r_w$, not zone-wide gradients.
    This explains the patch formation observed at scale S3--S4 (V94) without additional
    assumptions.
\end{itemize}

The empirical relay threshold $N_\mathrm{crit} = L/(5\,r_w)$ (Papers 35--37) acquires
a causal interpretation under this framing:
reliable zone-mean estimation requires $\geq 5$ independent causal neighbourhoods to average
over, providing adequate signal-to-noise on the zone-mean fieldM vector.
The constant 5 is a causal averaging threshold, not a geometric constant.

\textbf{Hierarchical relay (structure-about-structure).}
The finite causal horizon makes hierarchical relay a structural necessity, not an engineering
choice: a single region cannot produce coherent structure at scales exceeding $r_w$ because
no wave event can cross the horizon.
Trans-horizon coherence requires aggregation across causal neighbourhoods ---
which is exactly what relay coupling provides.
A distal region reads zone-level summaries from a proximal region (spatial averages
over $\geq 5$ causal neighbourhoods), crosses the causal horizon by aggregation,
and uses that summary as the input signal for its own causal neighbourhoods.
Key open question: does a two-level hierarchy produce sg4 of the \emph{summary}
at the distal region ($\equiv$ ``what does the pattern mean'', not just
``where is the pattern'')?

\textbf{Graph substrates.}
VCML-G (graph substrate) generalises the lattice by replacing Euclidean adjacency with
graph-topological adjacency.
The causal reparametrisation extends naturally: wave fanout $\phi_w$ becomes the mean
number of nodes reached by a single wave event on the graph; birth fanout $\phi_b$
(which is degenerate on the lattice, where births are purely local) becomes non-trivial
when newborn cells can sample fieldM from distant hub nodes via long-range edges.
The causal depth ratio $\delta = W_\mathrm{zone}/r_w$ generalises to
$\delta_G = d_\mathrm{zone}/r_G$, where $d_\mathrm{zone}$ is zone diameter in graph hops
and $r_G$ is the effective causal horizon in hops.
Key questions:
(1) Can tuning $\phi_b$ (long-range birth fanout) decouple wave radius from birth radius,
increasing effective zone coherence without increasing perturbation cost?
(2) Small-world connectivity increases $\phi_w$ without proportionally increasing wave cost---
does this overcome the $N_\mathrm{active}^{-1/2}$ scaling of sg4?
(3) Do $\phi_w$ and $\phi_b$ act independently on sg4, or do they compose multiplicatively
(i.e., is there a combined causal coupling $\Phi_\mathrm{total} = \phi_w \cdot \phi_b$)?

\textbf{Sequence learning.}
VCML currently encodes spatial (zone) structure.
Can the three-timescale architecture encode temporal sequences if the ``zone'' dimension is
replaced by temporal position (e.g., via oscillatory gating)?
Connection to the cortex-hippocampus theta sequence literature.

\textbf{The 8D optimum.}
Full optimisation over $(C_1, C_2)$ jointly.
Does the optimal point correspond to a biologically observed parameter ratio?

\subsection{Physics connections (conjectural)}

The following connections are not results --- they are conjectures with specified gaps.
They are included here because they direct future work and because stating the gap
explicitly is the only protection against treating speculation as established theory.

\textbf{Conjecture 1: Universality of the three-timescale architecture.}
The VCML three-timescale structure (h: fast per-step / m: medium per-wave / F: slow
intergenerational) may be universal rather than specific to VCML cells.
The proposed claim is: \emph{any adaptive memory system must have at least three timescales
with viability gates between them}; single-timescale systems are provably non-adaptive
under mandatory turnover.
The same structure appears in neurons (membrane potential / synaptic trace / structural
change), cortex (activation / STP / LTP), and evolution (genotype / allele frequencies /
species morphology).

\textit{What this would add if true:} a theorem about adaptive memory in general.
The non-obvious load-bearing part is the gate as \emph{necessary} (not merely useful):
without a consolidation criterion between timescales, fast-timescale noise overwrites
slow-timescale structure.

\textit{Gap that must close first:}
(i) Formal definition of ``adaptive memory'' under turnover;
(ii) proof that single-timescale systems fail this definition;
(iii) proof that three timescales are sufficient.
This is a theoretical result, not an experiment.
Until (i)--(iii) are proven, this is a motivated intuition, not a claim.

% MERA/AdS-CFT connection removed from theory document.
% The c(ℓ) conjecture has no closed prediction yet and reaches toward quantum gravity
% in a way that invites dismissal before the gap is closed.
% Full treatment in: physics_intuitions.md (Everything-Is-Timescales, MERA entries).

\subsection{Shannon's missing piece}

Shannon's channel capacity theorem (1948) characterises the \emph{ceiling}:
the maximum rate at which information can be transmitted through a noisy channel.
It tells you how much information \emph{can} arrive.
It does not tell you which information \emph{does} arrive --- which signals are
selected from the channel output and written into long-term storage.

This is not a gap Shannon overlooked; it is a different problem.
Selection requires a criterion external to the channel itself.
Shannon's framework is agnostic about what the receiver does with incoming bits.

The viability gate provides exactly this missing piece.
It is a selection criterion operating on the channel output:
writes to fieldM are admitted only when the site has been causally quiescent for
$\geq$ SS steps.
The gate is not a noise filter in Shannon's sense (it does not reduce channel noise);
it is a \emph{survival criterion} that determines which signals are consolidated into
the slow-timescale structural manifold.

\begin{quote}
  Shannon built the ceiling: maximum transmissible information.\\
  VCSM builds the floor: minimum survival criterion for consolidation.\\
  Together they characterise the information content of fieldM.
\end{quote}

What passes the gate is not the highest-amplitude signal, nor the most frequent,
nor the most recent --- it is the signal that arrives during causal quiescence,
i.e., the signal associated with successful recovery from a perturbation event.
The gate selects for \emph{viability-compatible} information.
In biological terms: the organism learns what it survived, not what it encountered.

The gap between Shannon's ceiling and VCSM's floor is the compression:
fieldM does not record the channel's full throughput.
It records the subset that was zone-specific, temporally settled, and
associated with a viable transition.
Zone structure is the residue of this selection.

\subsection{The G\"{o}del parallel}

G\"{o}del's incompleteness theorem (1931) establishes that sufficiently powerful
formal systems contain true statements that cannot be proven within the system.
The standard framing emphasises the logical gap: truth exceeds provability.
The VCML framing emphasises the \emph{structural cause} of that gap.

A formal system has no turnover.
Axioms are permanent.
Theorems, once derived, are permanent.
The system accumulates structure without any mandatory replacement.
A self-referential statement (``this statement is unprovable'') can exist
stably within the system's expression language indefinitely,
without ever resolving into a settled truth value.
It is what we might call an \emph{ungated} state:
a statement that neither consolidates nor collapses.

VCML survives the analogous paradox by construction.
Any state that fails to settle into causal quiescence within the consolidation window
is never written to fieldM.
Mid-memory decays under $\gamma_m = 0.99$ per step;
without a streak of $\geq$ SS calm steps, the accumulator is erased.
Additionally, cells collapse: the physical substrate that holds the unsettled state
is replaced.
There is no persistent ungated state.

The self-referential case: a cell whose hidden state tracks the current fieldM state
of its neighbours is chasing a moving target.
The neighbourhood fieldM changes as other cells collapse and are reseeded.
The calm streak never reaches SS.
The write never fires.
The self-referential pattern is structurally excluded from fieldM --- not by a logical
rule, but by the dynamics of mandatory turnover.

\begin{quote}
  Unprovable $=$ ungated. Same category.\\
  Formal systems collapse on self-reference because nothing ever dies.\\
  VCML survives because thoughts are mortal.
\end{quote}

The G\"{o}del gap is not a bug in logic.
It is what any system looks like when it has no turnover and no consolidation criterion.
VCML is the complementary architecture: mandatory turnover forces resolution.
States that do not settle are erased; states that settle are consolidated.
The system cannot accumulate ungated content because it has no mechanism to hold it.
Thoughts that do not resolve into a stable fieldM write disappear with the cell that held them.

% ─────────────────────────────────────────────────────────────────────────────
\section{Conclusion}

Learning from failure requires the failure to end.

VCML/VCSM provides a complete mechanistic account of zone-structured memory via mandatory
individual turnover.
Its core claim --- death-gated inheritance is the \emph{maintenance} mechanism for zone
structure --- is counterintuitive but robustly supported across eight layers of analysis
(mechanism, parameter geometry, spatial PDE, dynamics, multi-region coupling, continual
learning, rule-family robustness, adversarial resistance) spanning 44 papers,
with external validation of the causal geometry conjecture in GPT-2 (Paper~45)
and experimental grounding of the $C_2$ causal coordinates (Paper~46).
Formation is environmentally driven and does not require any specific mechanism;
maintenance is what the viability-gated copy-forward loop provides.
Paper~44 closes the formation--maintenance distinction with a direct adversarial test:
the adaptive mechanism maintains positive cosine fidelity with original zone structure
under 2000 steps of sustained adversarial input, while the passive mechanism encodes
the adversarial pattern ($\Delta = 0.205$, $p < 0.001$).
Memory, not imprinting.

The substrate-agnostic invariant (refined by Paper~46):
structure persists because slow-timescale writes are gated to causal quiescence windows.
In adversarial environments, this timing protects the original encoding by letting
the contaminating perturbation signal decay before writing.
In benign environments, writes at perturbation onset capture peak zone-specific signal.
\emph{The gate is a selectivity mechanism, not a noise filter;
its optimal threshold is set by adversarial pressure and proximity to the
commitment epoch $t^*$, not by signal quality (Paper~47).}
This claim survives graph substrates, arbitrary hidden-state dimensionality, and any
perturbation that does not alter the temporal structure of the quiescence gate.

The architecture has four genuinely distinctive properties:
\begin{enumerate}[itemsep=3pt]
  \item \textbf{Persistence via turnover.} Zone structure is maintained \emph{by} the
        death--birth cycle, not despite it. Optimal turnover rate is strictly positive.
  \item \textbf{Non-equilibrium.} The system is formally non-gradient
        (curl $\neq 0$, machine-verified). Memory is not an energy minimum.
        It is a driven steady state maintained by continuous wave input and field decay.
  \item \textbf{Perturbation amplification.} Periodic perturbation (task switching,
        boundary mismatch) does not erode zone structure --- it amplifies it $5$--$10\times$
        by sustaining the copy-forward loop.
  \item \textbf{Attractor stability as a direction toward symbolic behaviour (conjecture).}
        The crystallised regime (low $\nu$, high persistence) produces stable fieldM attractors:
        the system returns reliably to them under perturbation.
        This is demonstrated.
        Whether those attractors support \emph{compositional} symbolic reasoning --- and whether
        the generalisation ``any sufficiently stable copy-forward system produces symbolic
        behaviour'' holds --- has not been tested.
        Compositional symbol use requires that stable attractors compose under operations;
        this has not been demonstrated in VCML and is an open experimental question.
        The directional claim --- that symbolic capacity is a property of attractor dynamics,
        not of discrete logic gates --- is supported by the attractor evidence but is not
        established by it.
\end{enumerate}

These four properties are not incidental.
They are what any system must have if memory is to remain adaptive under environmental
change while retaining compositional structure, and they jointly define the architecture's
computational niche.

%──────────────────────────────────────────────────────────────────────────────
\section{VCSM as a Universality Class (Hypothesis)}

\subsection{The derivation}

Papers~52--56 establish a derivation structure: the VCSM architecture is not
engineered but forced. Given three inputs:
\begin{enumerate}
  \item \textbf{Primitive}: signed deviation from self-baseline
    ($\mathrm{hid} - \mathrm{baseline\_h}$, not raw amplitude, not global error)
  \item \textbf{Constraint}: carriers die (mandatory turnover)
  \item \textbf{Requirement}: distinctions must persist across carrier death
\end{enumerate}
each of the seven architectural components follows necessarily:

\begin{center}
\begin{tabular}{ll}
\toprule
Component & Why forced \\
\midrule
\texttt{baseline\_h}        & Primitive requires a reference state \\
\texttt{hid - baseline\_h}  & Primitive requires measuring departure \\
\texttt{mid\_mem}           & Fast events must couple to slow consolidation \\
\texttt{MID\_DECAY}         & Prevents accumulation of stale signal \\
Viability gate              & Prevents unstable cells from corrupting fieldM \\
\texttt{fieldM}             & Cells die; store must outlast the individual \\
Birth seeding               & Without inheritance, death resets structure \\
\bottomrule
\end{tabular}
\end{center}

The six experimental papers (Paper 43 onward) repeatedly confirm this: removing
any component produces exactly the failure mode predicted by the component's forced role.

\subsection{The universality-class claim}

The four conditions that force this architecture are:
\begin{enumerate}
  \item distinguishable states ($A \neq B$)
  \item mortal carriers
  \item persistence demand
  \item cross-generational transfer channel
\end{enumerate}
These are also the four conditions for Darwinian evolution: variation, mortality,
heredity, selection. VCSM may be Darwinian dynamics at the unit level --- the same
constraint structure operating inside a single substrate rather than at the population level.

\paragraph{Cross-substrate evidence.}
Immune memory and cultural transmission each map cleanly to all seven components
under completely different microscopic physics:

\begin{center}
\begin{tabular}{lll}
\toprule
VCSM component & Immune memory & Cultural transmission \\
\midrule
\texttt{baseline\_h}         & Naive repertoire         & Default norms \\
contrast signal              & Clonal expansion         & Marked deviation \\
\texttt{mid\_mem}            & Germinal centre          & Rehearsal / working memory \\
consolidation gate           & Affinity selection       & Social reinforcement \\
\texttt{fieldM}              & Memory B cells           & Institutional record \\
birth seeding                & Clonal inheritance       & Pedagogy \\
\bottomrule
\end{tabular}
\end{center}

Same forced architecture, different physics. That is the empirical signature of a
universality class.

\paragraph{Connection to physics.}
$C_\text{order}$ (Paper~55) is a scalar with monotone positive drift in the identity
regime --- a Lyapunov function. Lyapunov functions are the non-equilibrium analog of
Hamiltonians. The open question is whether $C_\text{order}$ can be derived from a
variational principle, as least action generates Newtonian mechanics. If yes, VCSM has
an action principle and the universality claim becomes a structural physics claim.

\paragraph{Falsification test.}
Find a system satisfying the four conditions that does \emph{not} produce the
seven-component architecture. If such a system exists, a constraint is missing or the
claim is wrong. If none can be found after systematic search across substrates, the
architecture is minimal and the universality class is real.

\paragraph{Status.}
This is a hypothesis, not a proven claim. The derivation is complete at the constraint
level; the universality requires showing the same architecture emerging from different
microscopic physics in at least two to three independent domains.

% ─────────────────────────────────────────────────────────────────────────────
\section{The True Primitive: Persistence of Distinguishability}
\label{sec:true_primitive}

\subsection{The full hierarchy (after ferromagnet counterexample)}

The ferromagnet counterexample (A\,$\neq$\,B + stable carrier = no VCSM needed)
revealed a branching point above VCSM: the super-primitive is not
``carrier turnover'' but \emph{persistence of distinguishability}.

\begin{center}
\begin{tabular}{ll}
\toprule
Level & Statement \\
\midrule
L0 & $A \neq B$ --- distinguishable states exist \\
L1 & Does $A \neq B$ still hold after time? \\
   & \textbf{``Persistence of distinguishability''} = super-primitive \\
L2a & Attractor persistence (local, passive): same field stays stable \\
    & ferromagnet, crystal, topological defect, Hopfield attractor \\
L2b & Reconstructive persistence (global, active): pattern rebuilt across reset \\
    & VCML fieldM, immune memory, evolution, culture \\
L3  & VCSM = minimal architecture of L2b \\
\bottomrule
\end{tabular}
\end{center}

\textbf{One-sentence formulation:}
\emph{There are two ways for a distinction to survive time:
stable attractor occupancy, or reconstruction across substrate change.
VCSM is the minimal architecture of the second.
A system can run both simultaneously.}

\subsection{The ferromagnet is absorbed, not refuted}

The ferromagnet is not outside the VCSM theory; it is the description of
VCML's \emph{local layer}.
Below the correlation length $r_\mathrm{corr}$, cells stabilise in
ordered patches (attractor persistence, ferromagnet-like).
Above $r_\mathrm{corr}$, zone identity is maintained only by the
copy-forward loop (reconstructive persistence).
V94 directly observed this: structure forms in patches, not zone-wide
gradients.
The patches are the attractor layer; zone-wide coherence is the
reconstructive layer.

\subsection{Coexistence of both layers simultaneously}

The two strategies are not alternatives but \emph{nested layers} operating
at different scales in the same system.
The correlation length $r_\mathrm{corr}$ is the boundary:
below it, attractor persistence suffices; above it, copy-forward reconstruction
is required.

In neuroscience, the same two-layer structure appears:
\begin{itemize}
  \item \emph{Local attractor layer:} cell assemblies (Hebb), Hopfield-style
    cortical attractors, working memory (persistent firing).
    The ``now'' of experience: moment-to-moment binding of a unified state.
  \item \emph{Global reconstructive layer:} hippocampal-to-cortical consolidation
    during sleep (batch copy-forward), long-term potentiation as fieldM,
    synaptic protein turnover as carrier mortality.
    The ``who'' of identity: narrative self persisting across sleep and
    neural turnover.
\end{itemize}

\subsection{Cross-substrate verification: Ising + VCSM-lite (Paper 60)}

A 2D Ising ferromagnet ($40 \times 40$, two spatial zones) augmented with a
VCSM-lite memory layer ($\phi$ field, causal gating, wave perturbations).
The system has two order parameters:
$m = |\langle s_i \rangle|$ (Ising, controlled by $T$) and
$S_\phi = |\bar\phi_{\rm z0}-\bar\phi_{\rm z1}|/\sigma_\phi$
(phi-order, controlled by $P_\mathrm{causal}$).

Key result: at $T=3.0$ ($T > T_c$, $m=0.07$, Ising disordered),
$S_\phi$ ranges from $0.51$ to $0.91$ as $P_\mathrm{causal}$
increases from $0.50$ to $1.00$.
Zone-specific memory order exists even when the underlying spin attractor has
been thermally destroyed.
The two control parameters are orthogonal.

\textbf{Causal temperature:} $(1-P_\mathrm{causal})$ plays the role of
thermal temperature for the purity-controlled transition.
It is an \emph{informational} noise amplitude, not a thermal one.
The universality class is ``purity-controlled non-equilibrium order,''
distinct from thermally controlled equilibrium order.

\textbf{Death as anti-crystallisation (Ising analogue):}
$S_\phi$ is lowest at $T=1.0$ (ordered spins, $m=1.00$): the wave
perturbation cannot move frozen spins, starving the reconstructive layer of
signal.
$S_\phi$ peaks near and above $T_c$.
Over-stable carriers (crystallised VCML, low-$T$ Ising) kill the reconstructive
layer, mirroring Papers~77--82--89.

\subsection{True Long-Range Order Confirmed: Extensive-Drive FSS (Paper 62)}

Paper~62 ran three tests; the decisive result is Phase~A.

\textbf{Phase A: extensive-drive FSS.}
Wave event rate scaled as WAVE\_EVERY$(L) = 25(40/L)^2$, holding wave density
per site constant.
Result: $|M| \propto L^{0}$ for $P_\mathrm{causal} \geq 0.10$ --- size-independent,
true long-range order in the thermodynamic limit.
The Binder cumulant $U_4$ is correctly ordered by $L$ for $P > 0.05$
(larger $L$ gives larger $U_4$, the ordered phase signature).
Critical point: $p_c \approx 0.02$--$0.05$.
The apparent transition near $P \approx 0.5$ in Paper~60 was a finite-size artifact
of intensive wave density.
Under extensive drive, even $\sim 3\%$ spatial causal purity sustains thermodynamic order.

\textbf{Phase B (Multi-Zone Relay):} phi-diffusion amplifies local zone separation
($\times 2.3$ at $D_\phi=0.10$); clean far-zone relay test deferred.

\textbf{Phase C (VCML Ablation):} null (minimal CML carrier too simple);
proper ablation deferred to Paper~65.

\textbf{Status of universality class.}
The causal-purity ordered phase is a genuine thermodynamic phase.
Paper~63 provides the precision measurement (below).

\subsection{Critical Exponents: Fine FSS (Paper 63)}

Fine scan $P_{\rm causal} \in [0.000, 0.080]$ in steps of $0.005$,
$L \in \{20,30,40,60,80\}$, 8 seeds, 12\,000 steps, extensive drive.

\textbf{Results:}
$p_c = 0.018 \pm 0.012$,
$\beta/\nu = 0.67$, $\nu = 0.97$, $\beta = 0.65$.
The susceptibility $\chi = N_{\rm zone} \cdot \mathrm{Var}(M)$ gives $\gamma/\nu \approx 0$
--- a definition artifact (zone-mean variance self-averages as $1/L^2$).
True $\gamma$ requires the spatial correlator $\sum_i \langle\phi_0\phi_i\rangle_c$.

\textbf{Key open question: is $p_c = 0^+$?}
For all $P_{\rm causal} \geq 0.005$, larger $L$ gives larger $U_4$ (ordered-phase signature),
suggesting the system may be ordered for \emph{any} non-zero causal purity.
If $p_c = 0^+$, the transition is infinite-order (Berezinskii--Kosterlitz--Thouless type),
not a standard second-order transition.
This question is answered in Paper~64 (below).

\textbf{Universality class comparison:}
$(\beta \approx 0.65,\ \nu \approx 1)$ matches no standard class
(mean field: $\beta=0.5,\ \nu=0.5$;
2D Ising: $\beta=0.125,\ \nu=1$;
directed percolation: $\beta=0.276,\ \nu_\perp=0.735$).
This is a candidate new non-equilibrium universality class.

\subsection{$p_c = 0^+$ Confirmed: BKT Test (Paper 64)}

Large-$L$ Binder crossing test:
$L \in \{40,60,80,100,120\}$,
$P_{\rm causal} \in \{0.000, 0.002, \ldots, 0.030\}$ (9 values),
12 seeds, 12\,000 steps, extensive-drive scaling.

\textbf{Phase~A result.}
$U_4(L=120) > U_4(L=40)$ for \emph{all} $P_{\rm causal} \geq 0.002$,
including the smallest non-zero value in the grid.
Binder crossings, where found, yield $p_c < 0.010$;
the majority of $L$-pairs produce no crossing at all.
Mean $p_c = 0.0079 \pm 0.0060$.

\textbf{Conclusion.}
The data are inconsistent with a finite $p_c$ of order $0.01$--$0.02$
and strongly consistent with $p_c = 0^+$:
the ordered phase exists for any positive causal purity.

\textbf{Phase~B result (spatial correlator).}
All correlation lengths $\xi = \mathrm{NaN}$;
the $\phi$-field amplitude $|\bar{M}| \sim 10^{-3}$ is too small
for reliable exponential fitting of $C(\Delta y)$.
The null result is consistent with $\xi \to \infty$ (algebraic order,
BKT-like) but does not confirm it.
The anomalous dimension $\eta(P_{\rm causal})$ and helicity modulus
were addressed in Paper~65 (below).

\textbf{Cosmological remark.}
$p_c = 0^+$ means ``causality $\Rightarrow$ order'' without threshold.
This is a structural complement to the causal set programme
(Bombelli, Lee, Meyer, Sorkin 1987):
if spacetime geometry emerges from causal partial ordering,
and thermodynamic $\phi$-order requires only nonzero causal purity,
then the existence of causality is sufficient for structure formation ---
dissolving rather than answering the initial-conditions fine-tuning problem.

\subsection{Probabilistic Wave Firing and $p_c=0^+$ Confirmation (Paper 67)}

Paper~67 introduces \emph{probabilistic wave firing} to correct the
\texttt{WAVE\_EVERY} integer-rounding artefact identified in Paper~66.
Instead of the deterministic rule $t \bmod w(L) = 0$, each step fires
a wave with probability $p_f(L)=\min(1,1/w_f(L))$ where
$w_f(L)=25(40/L)^2$ is the exact (non-integer) extensive-drive target.
For $L=160$: $p_f=0.640$, giving the correct $6400$ expected waves
per $10\,000$ steps (vs.\ $5000$ with the old deterministic WE=2).

\textbf{Phase A (discriminating test).}
With probabilistic firing at $L\in\{80,100,120,160\}$:
$U_4(160)>U_4(120)$ for $P\in\{0.010,0.020\}$ — the expected
ordered-phase monotonicity consistent with $p_c=0^+$.
At $P=0.005$ the margin is within 3-seed noise.

\textbf{Phase B (direct artefact quantification).}
At $L=160$, $P=0.020$: det-WE=2 gives $U_4=0.31$,
while probabilistic and det-WE=1 both give $U_4\approx0.50$.
The paper~66 reversal is directly attributed to the wave density deficit.

\textbf{New physics finding.}
At $P=0.005$: more wave density \emph{hurts} ordering (det-WE=1 gives
$U_4=0.097$, below det-WE=2's $U_4=0.283$).
At low causal purity, extra waves inject extra noise; the optimal wave
density is not ``maximum'' but depends on $P_\text{causal}$.

\textbf{Conclusion.} $p_c=0^+$ strongly supported for $P\ge0.010$.
Probabilistic firing is now the standard large-$L$ protocol.

\subsection{Discriminating Test, Direct $\phi$-Correlator, and $F_A$ Control (Paper 66)}

Three experiments resolved the two open questions from Paper~65.

\textbf{Wave-density rounding artefact (Phase~A).}
Phase~A repeats the large-$L$ Binder test at $L \in \{100,120,160\}$
with the same run parameters as Paper~65 and finds
$U_4(L=160) < U_4(L=120)$ at all $P > 0$.
The root cause is identified geometrically:
the formula \texttt{WAVE\_EVERY}$(L) = \max(1, \text{round}(25(40/L)^2))$
rounds to 2 for $L=160$ (continuous target: 1.5625),
giving $L=160$ only $78\%$ of the intended wave exposure per site
(15.2 hits vs.\ 18.0 for $L=120$).
This systematic under-driving, not a genuine phase boundary, explains the reversal.
$p_c = 0^+$ remains the favoured hypothesis; the direct confirmation requires
\texttt{WAVE\_EVERY}$=1$ for $L=160$ (deferred to Paper~67 due to compute cost).

\textbf{Null $\phi$-field spatial correlator (Phase~B).}
At boosted $F_A = 0.50$, the $\phi$-field amplitude reaches $|\phi| \approx 0.013$
(more than $10\times$ Papers~64--65).
Nevertheless, $G_\phi(r)$ is undetermined at all $(L,P)$:
the normalised correlator falls below the 2\% fit threshold at $r=1$.
The $\phi$-field has no spatial structure within a zone beyond 1 lattice spacing.
\emph{Ordering is entirely in the zone-mean $M$, not the spatial texture.}
The transition is a zone-mean symmetry-breaking, not BKT algebraic quasi-LRO.

\textbf{$F_A$ is an anti-ordering parameter (Phase~C).}
Sweeping $F_A \in \{0.10, 0.30, 0.50, 0.80\}$ at $L=80$, $P=0.020$ shows
$U_4$ decreasing monotonically with $F_A$ (from 0.31 to 0.06).
High $F_A$ causes $\phi$ to track the noisy mid-memory aggressively,
increasing variance and reducing zone separation.
$|\phi| \propto F_A^\alpha$ with $\alpha > 1$ (super-linear).
The optimal zone ordering is at low $F_A$, consistent with the
death-as-anti-crystallisation principle.

\subsection{Three Probes: Spin Correlator, Extended Binder, Ablation (Paper 65)}

Three further experiments tested the BKT hypothesis and identified
the load-bearing VCSM components.

\textbf{Spin correlator (Phase~A1).}
The Ising spin-spin autocorrelation $G_s(r)$ at $T=3.0$ is exponential
with $\xi \approx 1.1$ lattice spacings, independent of $P_{\rm causal}$.
The $\phi$-field ordering does not couple back to the Ising spin substrate;
the spin correlator is the wrong proxy.
The $\phi$-field stiffness $K_\phi \propto L$ is a wave-geometry artefact,
not a BKT stiffness.

\textbf{Extended Binder at $L=160$ (Phase~A2).}
$U_4(L=160) < U_4(L=120)$ for all $P_{\rm causal} \leq 0.005$,
reversing Paper~64's $p_c = 0^+$ conclusion.
Status: either $p_c$ is finite in $[0.005, 0.020]$,
or $L=160$ requires $>10\,000$ steps to equilibrate (WAVE\_EVERY rounds to 2
vs.\ the true extensive-drive value of 1.56).
The $p_c = 0^+$ claim is downgraded to an open question.

\textbf{VCSM ablation (Phase~C).}
NoGate (SS$=0$): $U_4 = 0.090$ (kills ordering).
The viability gate is the primary noise filter and is load-bearing.
Crystallized (FIELD\_DECAY$=1$): $U_4 = 0.262$ (highest), but via the
cockroach artefact --- $\phi$ accumulates without forgetting, locking into
an arbitrary early state.
This confirms that FIELD\_DECAY is a necessary anti-crystallisation mechanism.

\subsection{Finite-Size Scaling and Geometric Dilution (Paper 61)}

We attempted a full FSS analysis of the causal-purity transition in Ising + VCSM-lite,
measuring $|M|$, $\chi$, $U_4$, $S_\phi$ across $L \in \{20,30,40,60\}$
and $P_\mathrm{causal} \in [0,0.65]$.

\textbf{Main result.}
The order parameter scales as $|M| \propto L^{-2}$ at \emph{all} values of
$P_\mathrm{causal}$ (slope $-2.02 \pm 0.03$, independent of purity).
The Binder cumulant $U_4$ shows no $L$-dependent crossing.

\textbf{Interpretation: geometric dilution.}
Each wave event perturbs a fixed area $\approx \pi r_w^2$ while the zone
contains $L^2/2$ sites.
The signal density per site is $\propto L^{-2}$, independent of the physics.
This is not critical scaling; it is a fundamental constraint of the current
``intensive-drive'' protocol.

\textbf{Correct thermodynamic limit.}
For a true second-order transition, wave density must be extensive:
wave events $\propto L^2$ (i.e., WAVE\_EVERY $= O(1)$).
Under this scaling, $|M|$ becomes size-independent in the ordered phase
and standard FSS applies.

\textbf{What the Binder cumulant does confirm.}
Despite the absence of a crossing, $U_4$ transitions from $\approx 0$ (Gaussian)
to $\approx 2/3$ (bimodal) near $P_\mathrm{causal} \approx 0.05$--$0.10$,
independently of $L$.
The transition from Gaussian to bimodal phi-fluctuations is real.

\textbf{Status.}
The universality class claim stands as a hypothesis.
Paper~62 will run FSS under extensive-drive scaling to extract $\beta$, $\nu$, $\gamma$.

\bigskip
\noindent\textbf{Full paper index:} \url{https://github.com/AccidentalGenius101}\\
\noindent\textbf{Theory reference (Markdown):} \texttt{VCML\_THEORY.md} (Papers 0--33),

\subsection{Order Parameter Exponent, Functional Form, and Manna Falsification (Paper 68)}

Paper~68 conducts three simultaneous tests at the $p_c=0^+$ fixed point confirmed
in Paper~67.

\textbf{Phase A} (fine $P$ scan at $L=160$, $P\in[0.001,0.050]$, 6 seeds, 10\,000 steps)
discriminates among three functional forms for $|M(P)|$ as $P\to0^+$:
power law $|M|\sim P^\beta$ ($R^2=0.757$, best fit), BKT essential singularity
$|M|\sim e^{-A/\sqrt{P}}$ ($R^2=0.501$), and ordinary essential singularity
$|M|\sim e^{-A/P}$ ($R^2=0.307$).
The apparent $\beta_\mathrm{fit}=0.406$ is unreliable: for $P<0.010$ the Binder
cumulant $U_4<0.1$, indicating the system at $L=160$ is in the finite-size
disordered regime where $|M|$ reflects fluctuations rather than the genuine
order parameter.
Paper~63's FSS value $\beta=0.65\pm0.05$ remains the best estimate.

\textbf{Phase B} (FSS at fine $P$, $L\in\{80,100,120,160\}$)
shows scaling slopes $d\ln|M|/d\ln L$ evolving from $-1.21$ at $P=0.003$
(disorder limit $L^{-1}$) to $-0.53$ at $P=0.020$ (approaching the critical
value $-\beta/\nu=-0.670$).
This disorder-to-order crossover as $P$ increases is the expected signature of
a $p_c=0^+$ transition where there is no lower threshold.

\textbf{Phase C} (Manna conservation-law diagnostic, $L=120$, $P=0.010$) tests
a theoretically motivated universality class assignment.
The Manna/stochastic-sandpile class in 2D has $\beta\approx0.639$,
$\nu\approx0.799$, and a structurally appealing analogy: the copy-forward loop
maps onto conserved stochastic redistribution, $Z_2$ zone symmetry maps onto
Manna's local conservation law.
The $\beta$ proximity ($0.65$ vs $0.639$) is genuine.
However, Manna class requires an \emph{exactly conserved} activity variable
\citep{Vespignani1998}.
The diagnostic measures $\phi$-creation flux $\mathcal{F}_\mathrm{cr}=0.256$
(consolidation writes) and decay flux $\mathcal{F}_\mathrm{dc}=0.079$
(FIELD\_DECAY removal), yielding ratio $r=3.23\gg1$.
\textbf{The Manna class is falsified.}
VCSM belongs to a distinct universality class with directed dynamics,
$Z_2$ symmetry, and long-range spatial coupling via the wave field ($r_w=5$,
not nearest-neighbour).
The latter is absent from both Manna and directed percolation and is the
most likely source of the exponent displacement.

\textbf{Open question for Paper 69}: measure $\beta$ and $\nu$ at $r_w=1$
(nearest-neighbour wave coupling) to test whether the exponents converge
to a known class in the local-coupling limit.

\subsection{Wave-Range Dependence and Minimum Ordering Threshold (Paper 69)}

Paper~69 tests the conjecture from Paper~68 that the long wave range ($r_w=5$)
is responsible for displacing VCSM from known universality classes.

\textbf{Phase A} sweeps $r_w\in\{1,\ldots,8\}$ at fixed $L=80$, $P=0.020$.
The order parameter amplitude scales as $|M|\sim A(r_w)^{0.67}$, where
$A(r_w)=2r_w(r_w+1)+1$ is the Manhattan wave footprint.
$U_4$ rises from $0.15$ at $r_w=1$ to $0.34$ at $r_w=8$: ordering is
proportional to wave coverage, not just wave count.

\textbf{Phase B} attempts a matched-coverage FSS comparison between $r_w=1$ and $r_w=5$.
The $r_w=5$ FSS slopes evolve from $-0.88$ (disordered, small $P$) to $-0.21$
(ordered, $P=0.050$), consistent with previous papers.
The $r_w=1$ FSS slopes are $\approx-1.7$ at all $P$ --- steeper than the
disordered noise floor ($-1.0$) --- because the matched-coverage protocol breaks
down: the firing probability is capped at $p_\mathrm{wave}=1.0$ for all $L\geq40$
with $r_w=1$, so larger systems actually receive proportionally fewer hits per site.

\textbf{Phase C} (fine $P$ scan, $r_w=1$, $L=160$) shows $U_4<0.18$ for $P\leq0.015$
and an apparent $\beta_\mathrm{fit}=0.113$ ($R^2=0.63$), consistent with neither
DP ($\beta\approx0.584$) nor Paper~63 ($\beta=0.65$).
However, the system is not in the asymptotic ordered phase, so this measurement
is unreliable.

\textbf{Key finding}: The ordering transition in VCML requires a minimum effective
wave spatial range.
At $r_w=1$, the ordered phase is inaccessible at $P\leq0.050$ for $L=160$.
The threshold conjecture is $r_w\gtrsim\text{zone\_width}/20$.
The DP hypothesis from Paper~68 therefore cannot be tested cleanly with $r_w=1$
at the current system sizes; a protocol firing multiple waves per step would be
required.

\textbf{Paper~70}: joint $(r_w, P)$ scan at fixed $L$ to map the ordering threshold
$r_w^*(P)$ and determine whether it scales as $r_w^*\sim L^{1/2}$ or $L^1$.

\subsection{Threshold Map, Zone-Bleeding Anomaly, and RG Dimension of the Wave Operator (Paper 70)}

Paper~70 maps the ordering threshold $r_w^*(P)$ jointly across the $(r_w, P)$ plane,
resolves whether the Paper~69 minimum-range threshold is a finite-size artefact or a
thermodynamic property, and extracts the RG scaling dimension of the wave operator.

\textbf{Phase A} (2D phase map, $L=80$, Protocol A, $r_w\in\{1,2,3,4,5,6,8,10,12\}$,
$P\in\{0.005,\ldots,0.200\}$, 6 seeds, 8k steps):
$U_4>0.30$ (ordered) is achieved by $r_w=5$ at $P\geq0.020$ and by $r_w=3$ at $P\geq0.050$.
$r_w=1$ reaches $U_4>0.30$ only at $P\geq0.100$.
The threshold curve extracted from linear interpolation gives $r_w^*(P)\sim P^{-0.315}$
($R^2=0.77$) under constant-rate Protocol A.
A notable anomaly: $r_w=2$ is \emph{worse} than $r_w=1$ at several $P$ values.
The cause is zone bleeding: $r_w=2$ waves straddle the zone boundary and deposit
causal signal in both zones simultaneously, causing the zone-mean contrast $M$ to
cancel rather than amplify.

\textbf{Phase B} ($L$-dependence at $P=0.020$, Protocol A, $r_w\in\{1,\ldots,8\}$,
$L\in\{40,60,80,100,120\}$):
The threshold $r_w^*(L)$ \emph{decreases} with system size, with a fitted exponent
$\gamma\approx-1.08$ ($r_w^*\sim L^{-1.08}$).
This is opposite to the Paper~69 conjecture $r_w^*\sim L^{1/2}$.
The physical interpretation: as $L$ grows, the background fluctuation $\sigma_M\sim 1/L$
shrinks faster than the wave signal grows, so a smaller relative $r_w$ suffices.
\textbf{The Paper~69 minimum-range threshold is therefore a finite-size artefact.}
Any $r_w>0$ produces ordering in the thermodynamic limit $L\to\infty$.

\textbf{Phase C} ($r_w^*(P)$ under matched-coverage Protocol B, $L=80$,
$r_w\in\{1,\ldots,10\}$, $P\in\{0.005,\ldots,0.100\}$):
$r_w^*(P)\sim P^{-0.527}$ ($R^2=0.82$), very close to $P^{-1/2}$.
This confirms the signal-to-noise prediction: the zone-mean perturbation per wave
scales as $A(r_w)/N_\mathrm{zone}\sim r_w^2/L^2$, while fluctuations scale as $1/\sqrt{P}$,
giving threshold $r_w^*\sim P^{-1/2}$.

\textbf{RG dimension of the wave operator.}
The threshold relation $r_w^*\sim P^{-1/2}$ implies that the composite variable
\begin{equation}
  \xi \equiv r_w \cdot P^{1/2}
\end{equation}
is scale-invariant at the ordering threshold.
Under a spatial rescaling $\ell\to b\ell$ (RG step), $r_w\to r_w/b$ and $P\to Pb^{d_J}$
where $d_J$ is the scaling dimension of the causal-purity current $J$.
Scale-invariance of $\xi$ requires $d_J/2 = 1$, so $d_J=2$ and the wave operator
$W\sim r_w\sqrt{P}$ has RG eigenvalue (scaling dimension) $d_\mathrm{wave}=1/2$.
This inserts a non-local smeared interaction into the VCML field action:
\begin{equation}
  S[\phi] \supset g\int d^2x\,d^2y\; K_\sigma(x-y)\,J(x)\,\phi(y),
\end{equation}
where $K_\sigma$ is the wave kernel (Manhattan ball of radius $r_w$).
The operator $K_\sigma J$ has scaling dimension $d_J-d_\mathrm{wave}=3/2$,
making it a \emph{relevant} coupling that drives the ordering transition.

\textbf{Key findings}: (1) $r_w^*(P)\sim P^{-1/2}$ confirmed under matched-coverage protocol.
(2) $r_w^*$ decreases with $L$: Paper~69 threshold is finite-size only; any $r_w>0$ orders
thermodynamically. (3) Zone-bleeding anomaly at $r_w=2$ identified and explained.
(4) Wave operator scaling dimension $d_\mathrm{wave}=1/2$ extracted from threshold scaling.
(5) Causal-purity current dimension $d_J=2$, making the wave coupling relevant in the RG sense.

\textbf{Paper~71}: measure the full $\beta$-function by tracking $\beta$ and $\nu$
as a function of the composite variable $\xi/\xi^*$ (approach to the critical manifold).

\subsection{Iso-$\xi$ Universality, $\beta$ Extraction, and FSS Slopes (Paper 71)}

Paper~71 tests whether $\xi=r_w\sqrt{P}$ is universal and extracts $\beta$, $\nu$.

\textbf{Phase A} (iso-$\xi$ test, $L=80$, Protocol~A):
The test is \emph{negative}: at fixed $\xi$, $U_4$ is strictly ordered by $r_w$
(smaller $r_w$ wins) with spread $\Delta U_4 = 0.20$--$0.39$.
The root cause is that $A(r_w)\cdot P$ varies $42\%$ from $r_w=2$ to $r_w=8$
at fixed $\xi$ under Protocol~A (constant wave rate), because
$A(r_w) = 2r_w^2+2r_w+1$ grows faster than $r_w^2$.
Iso-$\xi$ universality is a Protocol~B property (equalized noise density).

\textbf{Phase B} ($\beta$ extraction, $r_w=5$, $L=80$, Protocol~A):
$|M|\sim(\xi-\xi^*)^\beta$ gives $\beta=0.628$, $R^2=0.922$ at $\xi^*=0.50$,
consistent with Paper~63 ($\beta=0.65$).
$r_w=8$ gives a shifted $\xi^*=0.65$ (amplitude deficit under Protocol~A)
and $\beta=0.494$; the $r_w=5$ estimate is taken as primary.

\textbf{Phase C} (FSS, $r_w=5$, $L\in\{40,60,80,100,120\}$):
Slopes $d\ln|M|/d\ln L$ fall monotonically from $-0.929$ ($\xi=0.500$) to
$-0.185$ ($\xi=1.118$).
The slope at $\xi=0.707$ is $-0.635$, consistent with $-\beta/\nu=-0.67$ and
locating $\xi^*\approx0.65$--$0.71$ for $r_w=5$ Protocol~A.

\textbf{Summary}: $\beta/\nu\approx0.64$, $\beta\approx0.63$, $\nu\approx0.98$,
indirect anomalous dimension $\eta=2\beta/\nu\approx1.28$ ($d=2$).
The exponent triple $(\beta,\nu,\eta)\approx(0.63,0.98,1.28)$ matches no known
2D non-equilibrium universality class; the large $\eta\approx1.28$ (vs
$\eta\approx0.23$ for Directed Percolation) is the fingerprint of the
non-local smeared wave coupling.
$\xi$ is a \emph{coherence} variable, not an amplitude variable:
it measures spatial efficiency of zone-mean driving per unit causal signal
and is universal only when noise density is equalized (Protocol~B).

\textbf{Paper~72}: direct measurement of $\eta$ from the two-point fieldM
correlator $G(r)\sim r^{-(d-2+\eta)}$ and check of scaling relation
$\beta=\nu\eta/2$ (in $d=2$).

\subsection{Two-Point Correlator Null Result: Zone-Mean Transition (Paper 72)}

Paper~72 attempts to measure $\eta$ directly from
$G(r)=\langle\varphi(x)\varphi(x{+}r)\rangle_c$ at criticality
($r_w{=}5$, $P{=}0.010$), across $L\in\{80,100,120,160\}$.

\paragraph{Main result.}
$G(r)\approx0$ at the cell level at all system sizes.
Correlator amplitudes are $O(10^{-9})$, at the statistical noise floor,
with $G(r{=}3)<0$ at three of four $L$ values.
No power-law $G(r)\sim r^{-\eta}$ is detectable.
This extends and confirms Paper~66's finding ($G(1)/G(0)<2\%$) to the critical
point and to larger $L$, ruling out a conventional local-field-theory description.

\paragraph{Physical interpretation.}
VCML ordering is a \textbf{zone-mean transition}: the order parameter
$M=\langle\varphi\rangle_{\mathrm{z}0}-\langle\varphi\rangle_{\mathrm{z}1}$
is a \emph{spatially averaged} (zero-dimensional) stochastic variable.
Cells within a zone fluctuate nearly independently; the coherent causal signal
(fraction $P_\mathrm{causal}$) is buried beneath per-cell noise but survives
in the zone average because errors cancel over ${\sim}L^2/2$ cells.
The anomalous dimension $\eta$ is therefore \emph{not defined} at the microscopic
level, and QFT ladder rung~6 is vacuous.

\paragraph{QFT ladder revision.}
Rungs~1--5 ($\varphi$, $\mathbb{Z}_2$, field equation, $d_\mathrm{wave}{=}1/2$,
independent $\beta$ and $\nu$) remain valid.
Rung~6 (direct $\eta$) is vacuous for a zero-dimensional order parameter.
The appropriate two-point function is the \emph{temporal} autocorrelator
$C(t)=\langle M(\tau)M(\tau{+}t)\rangle$, which defines the dynamic exponent~$z$.

\textbf{Summary}: $(\beta,\nu)\approx(0.63,0.98)$ describe the temporal
critical behaviour of the zone-mean $M(t)$.
$\eta$ is undefined.
\textbf{Paper~73}: temporal autocorrelator $C(t)$ and dynamic exponent~$z$;
check whether $z=2/\nu$ (diffusive) or $z\neq2/\nu$ (anomalous temporal scaling).

\subsection*{Paper 73: Dynamic Exponent $z$ from Temporal Autocorrelator --- Intrinsic VCSM Timescale Dominates}

Three phases targeted the dynamic exponent $z$ via the temporal autocorrelator
$C(t)=\langle|M(\tau)||M(\tau{+}t)|\rangle_c$ of the zone-mean order parameter.

\textbf{Phase A} ($r_w{=}5$, $P{=}0.010$, $L\in\{40,60,80,100,120\}$, 8 seeds, 30\,000 steps):
$\tau_L$ scatters between 214 and 742 steps with no monotone trend.
FSS fit $\tau_L\sim L^z$ gives $z_A=0.21$ with $R^2=0.04$ (noise-dominated).
No finite-size scaling of the relaxation time is observed.

\textbf{Phase B} ($r_w{=}5$, $L{=}80$, $P\in\{0.007,\ldots,0.050\}$, 8 seeds, 40\,000 steps):
$\tau_\mathrm{corr}\approx190$--$206$ steps across a 7-fold range in $P$ --- essentially flat.
Fit $\tau\sim P^{-z\nu}$ gives $z_B\approx0$ ($R^2=0.21$).
No critical divergence in $P\in[0.007,0.050]$.

\textbf{Phase C} ($r_w{=}5$, $P{=}0.010$, $L{=}120$, 8 seeds, 60\,000 steps):
Exponential fit: $\tau=647$, $R^2=0.602$.
Power-law fit: $C(t)\sim t^{-2.20}$, $R^2=0.800$ (power law wins by $\Delta R^2=0.198$).
The autocorrelator decays \emph{faster than exponential} in the window $t\in[100,4500]$.

\textbf{Key interpretation --- intrinsic VCSM timescale}:
The constant $\tau_\mathrm{VCSM}\approx200$ steps is not a critical divergence but an
\emph{architectural timescale} set by $\{\texttt{SS}=8,\,\texttt{FA}=0.30,\,\texttt{FIELD\_DECAY}=0.999\}$.
Two competing processes: slow exponential decay $\tau_\mathrm{slow}\approx1000$ steps
(FIELD\_DECAY) and fast discrete gate refresh $\tau_\mathrm{fast}\approx8$ steps (SS).
Their superposition produces the effective $\tau_\mathrm{VCSM}\approx200$ steps,
independent of $L$ and $P$ in the range tested.
The power-law $C(t)\sim t^{-2.2}$ is a \emph{mixture of exponentials}, not genuine scale-free dynamics.

\textbf{Why no divergence}:
The critical point is $P_c=0^+$; at $P\in[0.007,0.050]$ the system is deep in the ordered phase.
The critical timescale $\tau_\mathrm{crit}\sim P^{-z\nu}$ only dominates when
$\tau_\mathrm{crit}\gg\tau_\mathrm{VCSM}\approx200$, requiring
$P\ll P_\mathrm{crossover}\approx200^{-1/(z\nu)}\approx0.002$ (for $z=2$, $\nu=0.98$).
\textbf{Paper~74}: scan $P\in\{0.0001,0.0005,0.001,0.003\}$ approaching $P_c=0^+$
to observe $\tau\sim P^{-z\nu}$ once $\tau_\mathrm{crit}\gg200$.

\textbf{Current state of the VCML universality class (the ``Minecraft seed'')}:
\begin{align*}
  M   &\sim P^\beta,\quad\beta\approx0.63 \\
  L^* &\sim P^{-\nu},\quad\nu\approx0.98 \\
  \tau_\mathrm{VCSM} &\approx 200\ \text{steps (intrinsic gate timescale)} \\
  z   &= \text{?}\quad(\text{requires }P\to0^+)
\end{align*}
The pair $(\beta,\nu)\approx(0.63,0.98)$ is the static fingerprint of the universality class.
The dynamic exponent $z$ is the last digit of the seed.

%% ─────────────────────────────────────────────────────────────────────────────
\subsection*{Paper 74: Dynamic Exponent $z\approx0.48$ at $P\to0^+$ --- Sub-Diffusive Scaling}

With $P_\mathrm{crossover}\approx0.002$ established, Paper~74 scans $P\in[10^{-4},10^{-2}]$
using GPU-accelerated batch simulation (RTX~3060, \texttt{torch.compile}, 8 seeds batched).

\textbf{Phase A} (bridge region, $P\in\{0.001,\ldots,0.010\}$, $L{=}80$, 8 seeds, 12k steps):
$\tau_\mathrm{corr}$ is flat at $\approx762$--$1067$ steps --- a plateau well above
$\tau_\mathrm{VCSM}=200$.
This reveals that the phi-field order parameter has a longer intrinsic timescale
($\tau\approx1000$) than the Ising spin observable measured in Paper~73 ($\tau\approx200$).
The system is in the saturated ordered phase: $\xi(P)\gg L$ at $L=80$ for all
$P\geq0.001$.

\textbf{Phase B} (deep $P$, $P\in\{0.0001,0.0002,0.0005\}$, $L{=}80$, 8 seeds, 150k steps, GPU):
$\tau$ climbs substantially: $\tau(0.0005)=4812$, $\tau(0.0002)=9263$, $\tau(0.0001)=3772$.
The non-monotonicity at $P=0.0001$ is attributed to marginal equilibration or disordered-phase
entry; 150k steps at $\tau\approx4000$--$9000$ gives only 16--38 correlation times.

\textbf{Phase C} (FSS at $P{=}0.001$, $L\in\{40,60,80,100,120\}$, 8 seeds, 100k steps, GPU):
$\tau_L$ is anomalous (largest at $L=40$, non-monotone thereafter).
$U_4<0$ for $L\leq100$ (near-critical or disordered).
$z_C=-0.414$ ($R^2=0.133$) --- unphysical; FSS not applicable at $P=0.001$ for $L<160$.

\textbf{Global z fit (Phases A+B combined)}:
\begin{equation}
  \boxed{z \approx 0.48,\quad R^2=0.71,\quad \tau\sim P^{-z\nu},\quad\nu=0.98}
  \label{eq:p74z}
\end{equation}
This is \emph{sub-diffusive}: $z=0.48\ll z_\mathrm{ModelA}=2$ and $z<z_\mathrm{KPZ}=3/2$.
The result is consistent with ballistic wave-seeding providing a coherent (not random-walk)
channel for propagating the causal signal through the phi-field.

\textbf{Completed VCML universality class fingerprint}:
\begin{align*}
  M   &\sim P^\beta,\quad\beta\approx0.63 \\
  L^* &\sim P^{-\nu},\quad\nu\approx0.98 \\
  \tau &\sim P^{-z\nu},\quad z\approx0.48\ (\text{sub-diffusive}) \\
  \tau_\mathrm{floor} &\approx1000\ \text{steps (phi-field timescale)} \\
  P_c &= 0^+\ \text{(ordering at any }P>0)
\end{align*}
$z\approx0.48$ is the dynamic exponent of the VCML universality class.
Phase~C FSS at $P=0.001$ is noisy ($R^2<0.15$); Paper~75 addresses this limitation.

%%───────────────────────────────────────────────────────────────────────────
\subsection*{Paper 75: $z$ Probed at Deep $P$; $\gamma$ Null; FSS Blocking and Hyperscaling
Violation}

Paper~75 pushes into the deep-$P$ regime ($P<10^{-3}$) and attempts the first measurement
of the susceptibility exponent~$\gamma$ from $\chi=\mathrm{var}(M)$.
Both goals produce null results that are physically informative.

\paragraph{FSS blocking.}
The correlation length $\xi(P)\sim P^{-\nu}$ with $\nu=0.98$ grows rapidly as $P\to0^+$.
For $L=80$: $P_\mathrm{FSS}(80)\approx80^{-1.02}\approx10^{-2}$, so all Phase~B values
($P\in[3\times10^{-5},5\times10^{-4}]$) have $\xi(P)\gg L=80$.  In this FSS regime $\tau$
is controlled by finite-size effects, not critical divergence.  The fitted $z_B=-0.271$
($R^2=0.641$) is unphysical (negative sign).
True scaling at $P\sim10^{-4}$ would require $L\gg\xi(10^{-4})\approx6300$, which is
computationally inaccessible.  The Paper~74 scaling window $P\in[10^{-3},10^{-2}]$ is the
\emph{only} accessible regime; $z=0.48$ is confirmed by exclusion.

\paragraph{$\gamma$ null.}
$\chi=\mathrm{var}(M)\sim10^{-7}$ everywhere---disordered-phase noise.
The scaling $\chi\sim L^{-2}$ (Phase~C) is geometric dilution of the zone mean, not a
susceptibility divergence.  The cross-phase fit $\gamma_\mathrm{all}=0.382$ ($R^2=0.739$)
is contaminated by the different $L$ values of Phases~A ($L=160$) and~B ($L=80$).
A valid measurement requires the intensive susceptibility
$\chi_\mathrm{int}=L^2\,\mathrm{var}(M)$ in the \emph{ordered} phase ($U_4>0$, fixed $L$,
scanning $P$)---the Paper~76 protocol.

\paragraph{Hyperscaling violation.}
The VCML order parameter is a zone mean ($d_\mathrm{eff}=0$): Josephson's relation
$\nu d=2-\alpha$ gives $\alpha=2$, contradicting Rushbrooke's $\alpha\approx0.36$.
Standard QFT hyperscaling is inapplicable to the zone-mean system.

\paragraph{CUDAGraph speedup.}
A manual \texttt{torch.cuda.CUDAGraph} implementation (\texttt{vcml\_gpu\_v3.py}) gives
$229\,000$ steps/min at $(L=80,B=8)$ vs.\ ${\sim}75\,000$ for \texttt{mode='default'}:
$3.06\times$ speedup.  All future papers use~\texttt{v3}.

\textbf{Updated VCML universality class fingerprint}:
\begin{align*}
  M   &\sim P^\beta,\quad\beta\approx0.63 \\
  L^* &\sim P^{-\nu},\quad\nu\approx0.98 \\
  \tau &\sim P^{-z\nu},\quad z\approx0.48\
        (\text{sub-diffusive; valid for }P\in[10^{-3},10^{-2}]) \\
  \eta &: \text{undefined at cell level }(d_\mathrm{eff}=0\text{, zone-mean transition}) \\
  \gamma &: \approx0.71\text{ (indirect; Fisher scaling }\nu(2-\eta)\text{; direct \textbf{blocked} at }d_\mathrm{eff}=0) \\
  P_c &= 0^+
\end{align*}

\subsection*{Paper 76: $\chi_\mathrm{int}=L^2\,\mathrm{var}(M)$ --- Second $\gamma$ Null;
$d_\mathrm{eff}=0$ Kinematic Obstruction Confirmed}

Paper~76 applies the intensive susceptibility fix $\chi_\mathrm{int}\equiv L^2\,\mathrm{var}(M)$
to cancel geometric dilution, scanning three phases:
Phase~A ($L=80$, $P\in[0.003,0.050]$, 8 seeds, 300k steps),
Phase~B (FSS $P=0.010$, $L\in\{40\ldots160\}$, 200k steps),
and Phase~C ($L=120$, $P\in[0.002,0.030]$, 300k steps).

\paragraph{Result.}
$\chi_\mathrm{int}\approx0.0024$ flat across all phases, all $P$, all $L$.
Fitted exponents: $\gamma_A=0.009$ ($R^2=0.18$), $\gamma_C=-0.006$ ($R^2=0.18$),
FSS slope $\gamma/\nu=-0.126$ ($R^2=0.76$, negative = unphysical).

\paragraph{Kinematic explanation.}
Paper~72 confirmed $G_\phi(r)\approx0$ at all $r>0$ (no cell-level spatial correlations).
By the central limit theorem, $\mathrm{var}(M)\approx4\sigma_\phi^2/L^2$, so
$\chi_\mathrm{int}=L^2\,\mathrm{var}(M)\approx4\sigma_\phi^2=\mathrm{const}$.
The $L^2$ factor cancels the dilution but leaves a noise constant, not a
diverging quantity.  For $d_\mathrm{eff}=0$:
$\chi=\int d^dr\,G(r)=G(0)\cdot L^0=\mathrm{const}$---no thermodynamic divergence.
Direct $\gamma$ measurement is kinematically blocked.

\paragraph{Indirect estimate.}
Fisher scaling gives $\gamma=\nu(2-\eta)=0.98\times(2-1.28)=\mathbf{0.71}$,
with the caveat that Fisher's relation was derived for $d_\mathrm{eff}>0$.

\paragraph{Pure CUDA~C backend.}
\texttt{vcml\_gpu\_v4.py} and \texttt{vcml\_cuda/vcml\_step.cu} implement the full
VCML step in pure CUDA~C: 13 kernel types, 23 launches per step, cuRAND device API
for all random numbers, captured as a single CUDA graph.
Expected throughput: $350$--$500\,000$ steps/min at $(L=80,B=8)$, enabling $L\geq256$
FSS runs.

\textbf{Final VCML universality class fingerprint}:
\begin{align*}
  \beta   &\approx 0.628  &&\text{(direct: P63, P71)} \\
  \nu     &\approx 0.98   &&\text{(direct: P71)} \\
  z       &\approx 0.48   &&\text{(direct: P74; valid window }P\in[10^{-3},10^{-2}]\text{)} \\
  \eta    &\approx 1.28   &&\text{(indirect: }2\beta/\nu\text{; direct inaccessible)} \\
  \gamma  &\approx 0.71   &&\text{(indirect: }\nu(2-\eta)\text{; direct blocked at }d_\mathrm{eff}=0\text{)} \\
  P_c     &= 0^+          &&\text{(direct: P64, P67)}
\end{align*}
Universality class characterised by three directly measured exponents $(\beta,\nu,z)$
and the structural feature of $d_\mathrm{eff}=0$ zone-mean ordering.

\end{document}
